%=======================================================================
\chapter{Constitutive Equations}
%=======================================================================

The balance laws of Chapter 6 are true of every material. Mass is
conserved in steel and in honey; linear momentum balances in glass and
in air. That universality is what made them worth proving once and for
all, and it is also precisely why they cannot, by themselves, predict
anything at all. A law that does not distinguish steel from honey
cannot tell you what either of them will do.

The deficit can be counted exactly, and the count is the reason this
chapter exists. The unknown fields at a point of the body are the mass
density $\rho$, the motion $\bchi$ or equivalently the velocity $\bv$
(three components), the temperature $\theta$, the Cauchy stress $\bT$
(six components, the symmetry $\bT=\bT^{t}$ having been extracted in
Chapter 6 from the balance of moment of momentum), the heat flux $\bq$
(three), the internal energy density $\eps$ and the entropy density
$\eta$. That is
\[
   1+3+1+6+3+1+1=16
\]
scalar fields. Against them Chapter 6 supplies conservation of mass
(one equation), the balance of linear momentum (three) and the balance
of energy (one): five. The Clausius--Duhem inequality \eqr{6.187} adds
nothing to the tally, since an inequality selects among candidates but
determines none of them.

Eleven relations are therefore missing, and it is no accident that the
set $C'=\{\bT,\bq,\eps,\eta\}$ of \eqr{8.2} below has exactly eleven
components. Supply those eleven as functions of the rest and the system
closes: five equations left, five unknowns left, namely $\rho$, $\bchi$
and $\theta$. The eleven relations are the \emph{constitutive
equations}, and they are where the particular material enters the
theory. Everything before this chapter was true of all matter;
everything in it is true of some matter and false of the rest.

Such relations are idealizations. They are not deduced from anything
more fundamental within continuum mechanics, and they are judged solely
by whether the predictions they deliver agree with what is measured. For
any real substance that agreement is to be expected only within a
limited range of conditions, so a constitutive equation encodes both the
material and the circumstances: the same steel is elastic under a light
load, plastic under a heavy one, and molten at $1700$ degrees, and no
single relation covers all three. This book treats only the classical
cases, viscous and inviscid fluids and elastic solids, which are the
most fully developed and most widely used, and which are the
prerequisite for everything more elaborate.

What this chapter does \emph{not} do is write down a particular
constitutive equation. It asks the prior question: what may a
constitutive equation look like at all? Four answers are given to it,
and only the first is a free modelling choice.

\begin{itemize}\itemsep2pt
\item \emph{Locality}, Section~\ref{sec:ch8elastic}: the response at a
      material point is settled by what happens in an arbitrarily small
      neighbourhood of that point, and, in the class treated here, by
      the values of the deformation, the velocity and the temperature
      there to first order. This is an assumption, and a restrictive
      one.
\item \emph{Material symmetry}, Section~\ref{sec:ch8symm}: certain
      rearrangements of the reference neighbourhood are undetectable by
      any experiment, and the set of them is a group. Whether that group
      consists of rotations or of all volume-preserving maps is the
      difference between a solid and a fluid.
\item \emph{Noll's rule}, Section~\ref{sec:ch8noll}: the symmetry group
      belongs jointly to the material and to the reference
      configuration, and changing the reference conjugates the group.
      This is a theorem, not a choice.
\item \emph{Frame invariance}, Section~\ref{sec:ch8frame}: two observers
      must agree about the material, and the demand that they use the
      same response function restricts that function. This too is
      forced.
\end{itemize}

A fifth restriction, compatibility with the Clausius--Duhem inequality
of Section 6.7, is stated at the end of the chapter and exploited in the
two that follow.

%=======================================================================
\section{Elastic materials with heat conduction and viscosity}
\label{sec:ch8elastic}
%=======================================================================

The quantities for which constitutive equations are sought are the
present values of the stress $\bT$, the heat flux $\bq$, the energy
density $\eps$ and the entropy density $\eta$, all at a material point
$p\in B$ whose position in the reference configuration $\kappa$ is
$\bX$. We construct these relations from the standpoint of a single
observer $O$, and only in Section~\ref{sec:ch8frame} bring in a second
observer $O^{+}$ and ask what the two must agree about.

The order matters, and it is worth saying why we do not impose observer
agreement from the start. Material frame indifference is the
requirement that constitutive relations written by $O^{+}$ describe the
same material as those written by $O$. If they did not, then what a
substance \emph{is} would depend on who was watching it, and every
laboratory would have to publish its own elasticity of copper. Experience
says otherwise: material behaviour is intrinsic. Agreement between the
two observers is therefore not a further physical hypothesis but a
consistency condition, and consistency conditions are cleanest to apply
to a structure that already exists. So we build the structure first.

\subsection*{Locality, and the expansions it licenses}

The first restriction to suggest itself is \emph{locality}: the response
of the material at $p$ should depend on the values of the variables
constituting a process in a neighbourhood of $p$, that is, in a
neighbourhood of $\bX$, and on nothing further afield. This rules out
action at a distance: what the material does here is not directly
affected by what is happening a metre away, except through the
intervening material.

We narrow this further. The processes to which the material is deemed
sensitive are taken to consist of the current deformation
$\bchi_{\kappa}(\bX',t)$, the current material velocity
$\hat\bv(\bX',t)$ and the current temperature $\hat\theta(\bX',t)$ for
all $\bX'\in\kappa$ in an arbitrarily small neighbourhood of $\bX$. The
carets are the ones introduced in Remark~\ref{rem:decorations}: they
select the referential description, the function whose input is a place
in $\kappa$, which is the only description in which ``a neighbourhood of
the material point'' is a neighbourhood of a fixed set.

Assuming these to be smooth functions of position in $\kappa$, each may
be expanded about $\bX$. The expansions are ordinary first-order Taylor
expansions, and the coefficient of $\bX'-\bX$ in each is the referential
gradient of the quantity being expanded, evaluated at $\bX$:
\begin{equation}
\begin{aligned}
   \bchi_{\kappa}(\bX',t)&=\bchi_{\kappa}(\bX,t)+\bF(\bX,t)(\bX'-\bX)
     +o(|\bX'-\bX|),\\
   \hat\bv(\bX',t)&=\hat\bv(\bX,t)+\dot\bF(\bX,t)(\bX'-\bX)
     +o(|\bX'-\bX|),\\
   \hat\theta(\bX',t)&=\hat\theta(\bX,t)
     +\ip{\big(\Grad\hat\theta(\bX,t)\big)}{(\bX'-\bX)}
     +o(|\bX'-\bX|).
\end{aligned}
   \tag{8.1}\label{eq:8.1}
\end{equation}
Each of the three coefficients has already been computed in an earlier
chapter, and it is worth saying which is which, because the middle one
is the only one that is not a definition.

The first line is the statement that $\bF=\Grad\bchi_{\kappa}$, written
as an expansion; it is \eqr{3.102}, the local approximation of an
arbitrary deformation by a homogeneous one.

The second line has coefficient $\Grad\hat\bv$, and \eqr{4.2} showed
that $\Grad\hat\bv=\dot\bF$: differentiating $\chi_{i}$ with respect to
$X_{A}$ and then with respect to $t$ gives the same thing in either
order, so the referential gradient of the velocity is the material time
derivative of the deformation gradient. That interchange, and hence the
appearance of $\dot\bF$ rather than some new tensor, is the whole
content of the second line.

The third line is the definition of the gradient of a scalar field,
Theorem~\ref{thm:grad}, transplanted to $\kappa$.

The symbol $o(|\bX'-\bX|)$ is the Landau symbol of
Definition~\ref{def:landau}: a quantity which, divided by $|\bX'-\bX|$,
tends to zero as $\bX'\to\bX$. It is what makes ``for $|\bX'-\bX|$
sufficiently small these are dominated by the terms shown'' a precise
statement rather than a hope.

\begin{remark}[what the truncation is, and what it costs]
\label{rem:ch8trunc}
Nothing so far forces a stopping point. One could keep the next term in
each expansion, which for the first line would be the second gradient
$\Grad\bF$, and obtain a wider class of models. Truncating at first
order is a modelling decision and should be recognised as one.

The decision says: whatever the material at $p$ is sensitive to, it is
captured to sufficient accuracy by the value and the first derivative of
the deformation, the velocity and the temperature at $p$. If two
processes agree in those data at $p$, the material at $p$ cannot tell
them apart.

What it costs is every phenomenon carried by the higher gradients. A
material whose response involves $\Grad\bF$ has an intrinsic length
scale, since $\Grad\bF$ carries dimensions of one over length while
$\bF$ is dimensionless; such theories are used for surface tension, for
the width of a phase boundary and for the size effects seen in very
small specimens. None of that is available here. In exchange, the
theory becomes a theory of \emph{fields at a point}, and the whole
apparatus of Sections~\ref{sec:ch8symm} to \ref{sec:ch8frame} becomes
statements about tensors rather than about functionals.

A material whose response depends on the deformation only through $\bF$
is called \emph{simple}, following Noll. The present class is that of
simple materials with heat conduction and viscosity.
\end{remark}

\subsection*{The two lists}

Accordingly, if we mean to account for effects at this order, we assume
the constitutive response of the material, consisting of the set
\begin{equation}
   C'=\{\bT,\bq,\eps,\eta\},
   \tag{8.2}\label{eq:8.2}
\end{equation}
evaluated at $\bX$ and $t$, to be determined by the list
\begin{equation}
   L'=\{\bx,\bv,\theta,\bF,\dot\bF,\Grad\theta\},
   \tag{8.3}\label{eq:8.3}
\end{equation}
also evaluated at $\bX$ and $t$. Here $\bx=\bchi_{\kappa}(\bX,t)$,
$\bv=\hat\bv(\bX,t)$ and $\theta=\hat\theta(\bX,t)$: the list contains
the values at $\bX$ from \eqr{8.1} together with the three gradients
that accompany them.

It will prove more convenient to work not with $C'$ but with the
equivalent set
\begin{equation}
   C=\{\bT,\bq,\psi,\eta\},
   \tag{8.4}\label{eq:8.4}
\end{equation}
where $\psi$ is the Helmholtz free energy per unit mass, defined in
Chapter 6 by \eqr{6.186},
\[
   \psi=\eps-\theta\eta .
\]
The two sets carry the same information. Given $\theta$, which is in the
list $L'$ and therefore always available, the relation
$\psi=\eps-\theta\eta$ determines $\psi$ from $(\eps,\eta)$ and
$\eps=\psi+\theta\eta$ recovers $\eps$ from $(\psi,\eta)$; the map is a
bijection with $\eta$ and $\theta$ held fixed. Nothing is gained or lost
in the exchange, which is why \eqr{8.4} may replace \eqr{8.2} without
comment.

What is gained is convenience later. The local Clausius--Duhem
inequality was reduced in Chapter 6 to \eqr{6.187},
\[
   -\rho\big(\dot\psi+\eta\dot\theta\big)+\ip\bT\bL
   -\theta^{-1}\ip\bq{\grad\theta}\ \geq\ 0 ,
\]
in which $\psi$, not $\eps$, is the energy that appears. Since that
inequality is the instrument by which the chapters that follow will
extract restrictions on the constitutive functions, it pays to have the
constitutive list already written in the variable the inequality uses.
The exchange $\eps\mapsto\psi$ is a Legendre transformation trading the
entropy for the temperature as the independent thermal variable, and
$\theta$ is what appears in the list \eqr{8.3}, not $\eta$.

\subsection*{Two gradients of the temperature}

The list $L'$ contains $\Grad\theta$, a referential vector. The list we
shall actually use contains $\grad\theta$, a spatial one. They determine
each other, and the relation between them is worth deriving rather than
quoting, since the same two-line argument recurs three more times in
this chapter.

Let $\dl\bX$ be a material line element at $\bX$ and $\dl\bx=\bF\dl\bX$
its image, \eqr{2.34}. The increment of temperature between the two ends
of that element is one number, and each description computes it its own
way:
\begin{equation}
   \ip{\Grad\hat\theta}{\dl\bX}=\dl\theta
   =\ip{\grad\tilde\theta}{\dl\bx}
   =\ip{\grad\tilde\theta}{\bF\,\dl\bX},
   \tag{8.5}\label{eq:8.5}
\end{equation}
the first and third equalities being the definition of the gradient in
$\kappa$ and in $\kappa_{t}$ respectively, and the last the substitution
$\dl\bx=\bF\dl\bX$. Now move $\bF$ across the inner product using the
definition of the transpose, Definition~\ref{def:transpose},
$\ip\ba{\bF\bb}=\ip{\bF^{t}\ba}\bb$:
\[
   \ip{\grad\tilde\theta}{\bF\,\dl\bX}
   =\ip{\bF^{t}\grad\tilde\theta}{\dl\bX}.
\]
The two ends of \eqr{8.5} are therefore inner products of $\dl\bX$ with
two fixed vectors, and they agree:
\[
   \ip{\big(\Grad\hat\theta-\bF^{t}\grad\tilde\theta\big)}{\dl\bX}=0 .
\]
The element $\dl\bX$ was arbitrary, so Lemma~\ref{lem:cvec}, the
constant-vector argument of Chapter 5, applies and the bracket vanishes:
\begin{equation}
   \Grad\hat\theta=\bF^{t}\grad\tilde\theta .
   \tag{8.6}\label{eq:8.6}
\end{equation}
This is Theorem~\ref{thm:GradgradF} of Chapter 2 specialised to a scalar
field. Note that it is the arbitrariness of $\dl\bX$ that does the work,
exactly as Remark~\ref{rem:cancel} warned: from a single $\dl\bX$ one
could conclude nothing, since any vector perpendicular to it would
satisfy the display.

Since $\bF$ is invertible, \eqr{8.6} may be read either way,
$\grad\tilde\theta=\bF^{-t}\Grad\hat\theta$, so the two gradients
determine each other once $\bF$ is known. And $\bF$ is known: it is
already in the list. The same is true of the other exchange we are about
to make. By \eqr{4.3}, $\bL=\dot\bF\bF^{-1}$, so $\dot\bF=\bL\bF$ and
$\bL=\dot\bF\bF^{-1}$ pass between $\dot\bF$ and $\bL$ in either
direction, again using only $\bF$. Hence $L'$ is equivalent to the list
\begin{equation}
   L=\{\bx,\bv,\theta,\bF_{\kappa},\bL,\grad\theta\},
   \tag{8.7}\label{eq:8.7}
\end{equation}
where $\bL$ is the spatial velocity gradient of \eqr{4.3} and the
subscript $\kappa$ has been attached to the deformation gradient to
record the reference configuration relative to which it is computed.
None of the other entries of \eqr{8.7} depends on a reference
configuration: $\bx$, $\bv$, $\theta$, $\bL$ and $\grad\theta$ are all
spatial, and $\bL$ in particular was shown in Remark~\ref{rem:whyL} to
be indifferent to the choice of $\kappa$. The deformation gradient is
the sole entry that is not, and the whole of
Sections~\ref{sec:ch8symm} and \ref{sec:ch8noll} is the consequence of
that single fact.

\begin{remark}[three letters that look alike]\label{rem:ch8letters}
Equation \eqr{8.7} contains both $L$ and $\bL$, and they are unrelated.
The plain italic $L$ is the \emph{list}, a set of six entries; the
under-tilde $\bL$ is the spatial velocity gradient, one of them. The
book distinguishes them by boldface, we by the under-tilde, and the
distinction is not decorative: the sentence ``$\bL\in L$'' is true and
``$L\in L$'' is nonsense.

The same warning will be needed twice more. From \eqr{8.10} onward,
$F$ without an under-tilde denotes a generic member of the response set
$C$, so that $F$ may stand for the stress, the heat flux, the free
energy or the entropy, while $\bF$ remains the deformation gradient and
$F_{\kappa}$ is a \emph{function}, not a value. And $\kappa$ denotes the
reference configuration while $\bkappa$ denotes the map $p\mapsto\bX$
that produces it. In each pair the two objects sit in the same formula,
so reading the decoration is not optional.
\end{remark}

\subsection*{Why this particular list}

The constitutive assumption adopted here may be motivated by an analogy
with a spring and a damper. The force in a spring is fixed by its
present elongation; the force in a damper is fixed by the present rate
of elongation. A device made of the two in parallel responds to the pair
(elongation, rate of elongation) and to nothing else, in particular not
to how it arrived at that state. The list \eqr{8.7} is the continuum
version of that pair: $\bF_{\kappa}$ is the local elongation and $\bL$
its rate. So the class of materials studied here is the class that
exhibits both elasticity and viscosity, and we should expect the
constitutive functions to split, eventually, into a part depending on
$\bF$ that stores energy and a part depending on $\bL$ that dissipates
it. The chapters that follow carry that split out, Chapter 9 doing it
for viscous fluids first.

\begin{remark}[what this class cannot do]\label{rem:ch8viscoel}
The analogy also marks the boundary. A real polymer stretched and held
does not settle instantly to the force its present elongation would
dictate; the force decays over minutes or hours. Its stress at time $t$
depends on the whole history $\bF(s)$ for $s\leq t$, not on $\bF(t)$ and
$\dot\bF(t)$ alone. The response is then a \emph{functional} on a space
of histories rather than a function on a finite list, and the entire
apparatus of this chapter, which manipulates the arguments of a
function, has to be rebuilt.

That is the general theory of materials with memory, and it lies outside
the scope of this book. The class treated here is not general
viscoelasticity, and it should not be represented as such. It is,
however, enough for the two applications that follow, and it is the case
in which every idea of the subject can be seen without the analytic
machinery that histories require.
\end{remark}

%=======================================================================
\section{Change of reference configuration and material symmetry}
\label{sec:ch8symm}
%=======================================================================

The list $L$ of variables presumed to constitute the relevant process
depends on the choice of reference configuration, and depends on it
through exactly one entry, the deformation gradient. To see what that
dependence does to the theory, change the reference configuration from
$\kappa$ to $\mu$ (Figure~\ref{fig:ch8chref}).

\bookfig{1}\begin{figure}[htbp]\centering
\begin{tikzpicture}[scale=1.0,
   ar/.style={-{Stealth[length=2.4mm]},thick}]
% ---------------- the body B (top) ----------------
\draw[bodyline] plot[smooth cycle,tension=0.85] coordinates
  {(0,4.13) (0.64,3.93) (0.83,3.33) (0.34,2.91) (-0.43,3.03)
   (-0.72,3.58)};
\node[font=\small] at (-1.30,4.12) {$B$};
\fill (0,3.55) circle (1.4pt);
\node[font=\small,anchor=west] at (0.10,3.59) {$p$};
% ---------------- kappa (left) ----------------
\draw[configline] plot[smooth cycle,tension=0.85] coordinates
  {(-4.30,1.18) (-3.66,0.98) (-3.47,0.38) (-3.96,-0.04) (-4.73,0.08)
   (-5.02,0.63)};
\node[font=\small] at (-5.55,1.35) {$\kappa$};
\fill (-4.30,0.60) circle (1.4pt);
\node[font=\small,anchor=west] at (-4.20,0.64) {$\bX$};
% ---------------- mu (right) ----------------
\draw[configline] plot[smooth cycle,tension=0.85] coordinates
  {(4.30,1.18) (4.94,0.98) (5.13,0.38) (4.64,-0.04) (3.87,0.08)
   (3.58,0.63)};
\node[font=\small] at (5.65,1.35) {$\mu$};
\fill (4.30,0.60) circle (1.4pt);
\node[font=\small,anchor=west] at (4.40,0.64) {$\bY$};
% ---------------- kappa_t (bottom) ----------------
\draw[configline] plot[smooth cycle,tension=0.85] coordinates
  {(0,-1.87) (0.64,-2.07) (0.83,-2.67) (0.34,-3.09) (-0.43,-2.97)
   (-0.72,-2.42)};
\node[font=\small] at (1.35,-3.05) {$\kappa_{t}$};
\fill (0,-2.45) circle (1.4pt);
\node[font=\small,anchor=west] at (0.10,-2.41) {$\bx$};
% ---------------- the four maps ----------------
\draw[ar] (-0.90,3.20) to[bend right=14] (-3.45,1.25);
\node[font=\small,lbl] at (-2.95,2.75) {$\bX=\bkappa(p)$};
\draw[ar] (0.90,3.20) to[bend left=14] (3.45,1.25);
\node[font=\small,lbl] at (2.95,2.75) {$\bY=\vek{\mu}(p)$};
\draw[ar] (-3.95,-0.10) to[bend right=14] (-0.95,-2.05);
\node[font=\small,lbl] at (-3.45,-1.65) {$\bx=\bchi_{\kappa}(\bX,t)$};
\draw[ar] (3.95,-0.10) to[bend left=14] (0.95,-2.05);
\node[font=\small,lbl] at (3.45,-1.65) {$\bx=\bchi_{\mu}(\bY,t)$};
% ---------------- lambda, the induced map ----------------
\draw[ar,red!65!black,densely dashed] (-3.30,0.60) -- (3.42,0.60);
\node[font=\small,lbl,text=red!65!black] at (0,0.60)
  {$\bY=\vek{\lambda}(\bX)$};
\end{tikzpicture}
\caption{Change of reference configuration. The body $B$ is a set of
material points and is not a region of space; $\bkappa$ and
$\vek{\mu}$ are two ways of laying it out in $\hEt$, and
$\bchi_{\kappa}$, $\bchi_{\mu}$ are the two descriptions of one and the
same motion. Because both placements are one to one, the diagram can be
closed along the top: $\vek{\lambda}=\vek{\mu}\circ\bkappa^{-1}$ takes
$\kappa$ to $\mu$ without passing through $B$, and it is the gradient of
this map, and nothing about $B$ itself, that enters \eqr{8.8} and
\eqr{8.9}. The current configuration $\kappa_{t}$ is reached by two
routes and does not care which was taken; that is the content of
$\bF_{\kappa}=\bF_{\mu}\bR$.}
\label{fig:ch8chref}
\end{figure}

Because the body and any reference configuration are in one-to-one
correspondence, there is an invertible map $\vek{\mu}(p)$ taking $p\in
B$ to $\bY\in\mu$. And $B$ and $\kappa$ are in one-to-one
correspondence as well, so the same is true of $\kappa$ and $\mu$:
composing the inverse of one placement with the other gives a map from
$\kappa$ straight to $\mu$ that never mentions $B$. We assume it smooth,
and call it $\vek{\lambda}$:
\begin{equation}
   \bY=\vek{\lambda}(\bX)
   \quad\text{and}\quad
   \dl\bY=\bR(\bX)\dl\bX,
   \quad\text{where }\bR(\bX)=\Grad\vek{\lambda}.
   \tag{8.8}\label{eq:8.8}
\end{equation}
The second statement is the first differentiated:
$\vek{\lambda}$ is a differentiable vector field on $\kappa$, so by
Definition~\ref{def:diffv} its increment is
$\vek{\lambda}(\bX+\dl\bX)-\vek{\lambda}(\bX)
=(\Grad\vek{\lambda})\dl\bX+o(|\dl\bX|)$, and the tensor
$\Grad\vek{\lambda}$ carrying $\dl\bX$ to $\dl\bY$ is what \eqr{8.8}
calls $\bR$.

The deformation gradients relative to $\kappa$ and to $\mu$ are now
related. Take one material line element and write its image in
$\kappa_{t}$ in the two ways the diagram permits:
\begin{equation}
   \bF_{\kappa}\dl\bX=\dl\bx=\bF_{\mu}\dl\bY=\bF_{\mu}\bR\,\dl\bX;
   \qquad\text{hence}\qquad
   \bF_{\kappa}=\bF_{\mu}\bR .
   \tag{8.9}\label{eq:8.9}
\end{equation}
The word ``hence'' is doing a small job here and it should be watched.
What the first chain gives is
$\big(\bF_{\kappa}-\bF_{\mu}\bR\big)\dl\bX=\bze$ for the particular
$\dl\bX$ chosen; but $\dl\bX$ was any material element at $\bX$, so the
tensor in the bracket annihilates every vector of $\hVt_{\kappa}$, and a
tensor that sends every vector to $\bze$ is the zero tensor by
Definition~\ref{def:tensor}. That is what licenses cancelling $\dl\bX$,
and it is the same quantifier argument as in \eqr{8.6}.

Two consequences of \eqr{8.9} are worth extracting at once. Taking
determinants and using Theorem~\ref{thm:det},
\[
   \det\bF_{\kappa}=(\det\bF_{\mu})(\det\bR),
\]
and both deformation gradients have positive determinant because both
$\kappa$ and $\mu$ are occupiable configurations; hence $\det\bR>0$ and
$\bR$ is invertible. Second, $\bR$ does not depend on $t$: it is built
from two reference configurations, neither of which moves.

\begin{remark}[this $\bR$ is not the rotation of the polar
decomposition]\label{rem:ch8Rwarning}
The letter is unfortunate and the book flags it explicitly, so we do
too. The $\bR$ of \eqr{3.32}, $\bF=\bR\bU=\bV\bR$, is orthogonal and
maps $\hVt$ to $\Vt$, from the reference translation space to the
current one. The $\bR$ of \eqr{8.8} maps $\hVt_{\kappa}$ to
$\hVt_{\mu}$, from one \emph{reference} translation space to another,
and it is not orthogonal: nothing has been assumed about it beyond
invertibility and $\det\bR>0$.

The index test of Remark~\ref{rem:hatinch4} separates them at a glance.
The polar rotation carries $\be_{i}\otimes\hbE_{A}$, one Latin index and
one capital. The present $\bR$ carries $\hbE_{A}\otimes\hbE_{B}$, two
capitals, and that is exactly why it can stand on the right of
$\bF_{\mu}$ in \eqr{8.9} and produce a two-point tensor again.
Confusing the two would make $\bR^{t}\bR=\hI$ look like a triviality
instead of what it becomes in \eqr{8.19}, namely the defining property
of a solid.
\end{remark}

\subsection*{The response function, and why it must change with
$\kappa$}

Let $F\in C$ be an entry in the list of variables characterizing the
response of the material: $F$ may be $\bT$, or $\bq$, or $\psi$, or
$\eta$: whichever of the four is under discussion. Its value is
determined by the entries of $L$, which we have just seen depends on the
choice of reference. Choosing $\kappa$,
\begin{equation}
   F(p,t)=F_{\kappa}(\bx,\bv,\theta,\bF_{\kappa},\bL,\grad\theta;\bX),
   \tag{8.10}\label{eq:8.10}
\end{equation}
where the right-hand side is the constitutive function delivering the
values of $F$. The semicolon is not a typographical flourish: the
entries before it are the process variables, which change from instant
to instant, while $\bX$ after it identifies which material point is
being asked about and is a parameter, fixed once the point is chosen.
The dependence on $\bX$ is what allows a body to be inhomogeneous.

\begin{remark}[the value does not depend on $\kappa$; the function
must]\label{rem:ch8valfun}
This is the pivot of the whole section, and it is one sentence: $F(p,t)$
is a number, or a vector, or a tensor, that the material actually has,
whereas $F_{\kappa}$ is a rule we wrote down for computing it.

The rule takes $\bF_{\kappa}$ among its arguments, and $\bF_{\kappa}$
changes when $\kappa$ changes. The output does not change: the stress
at $p$ at time $t$ is whatever it is, and no bookkeeping we perform can
alter it. So if the input changes and the output does not, the rule must
change to compensate. That is the entire argument, and \eqr{8.12} is its
statement.

This is the same distinction as Remark~\ref{rem:decorations}, where one
velocity was carried by three different functions $\bv$, $\hat\bv$,
$\tilde\bv$ according to what was fed in. Here one response is carried
by a family of functions $F_{\kappa}$, one for each reference
configuration. Miss the distinction and \eqr{8.16}, the definition of a
symmetry transformation, reads as the false claim that $\bF$ and
$\bF\bR$ are the same tensor. They are not. The claim is that a
particular function does not distinguish them.
\end{remark}

Using \eqr{8.8} and \eqr{8.9} we may rewrite \eqr{8.10} in terms of
$\mu$. The recipe is mechanical: wherever $\bX$ appears, substitute
$\bX=\vek{\lambda}^{-1}(\bY)$, and wherever $\bF_{\kappa}$ appears,
substitute $\bF_{\kappa}=\bF_{\mu}\bR$ with $\bR$ evaluated at that same
$\bX$. Then
\begin{equation}
   F(p,t)=F_{\mu}(\bx,\bv,\theta,\bF_{\mu},\bL,\grad\theta;\bY),
   \tag{8.11}\label{eq:8.11}
\end{equation}
where
\begin{equation}
\begin{aligned}
   &F_{\mu}(\bx,\bv,\theta,\bF_{\mu},\bL,\grad\theta;\bY)\\
   &\qquad=F_{\kappa}\big(\bx,\bv,\theta,
     \bF_{\mu}\bR(\vek{\lambda}^{-1}(\bY)),\bL,\grad\theta;
     \vek{\lambda}^{-1}(\bY)\big),
\end{aligned}
   \tag{8.12}\label{eq:8.12}
\end{equation}
which produces the constitutive function pertaining to any choice of
reference configuration once the function pertaining to one is given.

Read \eqr{8.12} carefully, because it is a definition dressed as an
equation. The left-hand side is not known in advance; the right-hand
side is a formula built from the known $F_{\kappa}$, from
$\vek{\lambda}$ and from $\bR$, and \eqr{8.12} defines $F_{\mu}$ to be
that formula. What is being asserted, then, is not the equality itself
but two things that follow from it.

First, $F_{\mu}$ exists: a change of reference never destroys the
possibility of a constitutive description, it merely relabels it.
Second, and this is the substantive point, the \emph{list of arguments}
is unchanged. On the right-hand side of \eqr{8.12} the fourth slot is
occupied by a deformation gradient, the sixth by a spatial temperature
gradient, and so on, exactly as on the left. No new kind of variable had
to be introduced to make the transfer work. Hence the list of variables
deemed relevant to material response is the same for all reference
configurations, even though the function of them is not. Had this failed,
the constitutive assumption of Section~\ref{sec:ch8elastic} would have
been an assumption about $\kappa$ rather than about the material, and it
would have been worthless.

\subsection*{Localizing the change of reference}

Constitutive response is defined pointwise, so nothing is lost by
confining attention to changes of reference in a neighbourhood of one
material point $p_{0}$, rather than over the whole body
(Figure~\ref{fig:ch8local}). Let $\vek{\lambda}(\bX)$ be smooth and
invertible on an open neighbourhood $N_{\kappa}(p_{0})\subset\kappa$,
mapping it onto another open neighbourhood $N_{\mu}(p_{0})$, and let
$\bX_{0}=\bkappa(p_{0})$ be a \emph{pivot point} of the map:
\[
   \bY_{0}=\vek{\lambda}(\bX_{0})=\bX_{0}.
\]

\bookfig{2}\begin{figure}[htbp]\centering
\begin{tikzpicture}[scale=1.0,
   ar/.style={-{Stealth[length=2.4mm]},thick}]
% ------------- kappa, with the neighbourhood inside it -------------
\draw[configline] plot[smooth cycle,tension=0.85] coordinates
  {(-3.10,1.50) (-1.65,1.00) (-1.35,-0.40) (-2.45,-1.60) (-4.10,-1.35)
   (-4.75,0.05)};
\node[font=\small] at (-5.15,1.30) {$\kappa$};
\draw[guide] (-3.05,0.05) circle (0.66);
\fill (-3.05,0.05) circle (1.5pt);
\node[font=\small,anchor=south] at (-3.05,0.16) {$p_{0}$};
\node[font=\small] at (-3.05,2.40) {$N_{\kappa}(p_{0})$};
\draw[gray!70] (-3.05,2.15) .. controls (-2.85,1.55) .. (-2.70,0.62);
% ------------- the image neighbourhood -------------
\draw[guide] (3.05,0.35) circle (0.66);
\fill (3.05,0.35) circle (1.5pt);
\node[font=\small,anchor=south] at (3.05,0.46) {$p_{0}$};
\node[font=\small] at (3.05,2.40) {$N_{\mu}(p_{0})$};
\draw[gray!70] (3.05,2.15) .. controls (2.85,1.60) .. (2.70,0.97);
\node[font=\footnotesize,text=gray!80,align=left,anchor=west]
  at (3.95,-0.80) {$\mu$ need exist\\ only here};
% ------------- the map -------------
\draw[ar] (-2.30,0.30) to[bend left=12] (2.30,0.60);
\node[font=\small,lbl] at (0,1.20) {$\bY=\vek{\lambda}(\bX)$};
% ------------- the pivot -------------
\draw[thinvec,red!65!black] (-1.55,-2.30) -- (-3.20,-0.42);
\node[font=\small,text=red!65!black,anchor=north] at (-1.50,-2.35)
  {$\bX_{0}$};
\draw[thinvec,red!65!black] (1.60,-2.30) -- (2.92,-0.20);
\node[font=\small,text=red!65!black,anchor=north] at (1.65,-2.35)
  {$\bX_{0}$};
\node[font=\footnotesize,text=red!65!black,anchor=north] at (0,-3.05)
  {one and the same point: $\vek{\lambda}(\bX_{0})=\bX_{0}$};
\end{tikzpicture}
\caption{Local change of reference. Only a neighbourhood of the one
material point $p_{0}$ is relabelled, so $\vek{\lambda}$ need not, and
in general does not, extend to a map of the whole body; $\mu$ is a local
reference configuration and nothing more. The two arrows on the right
carry the same label because $\bX_{0}$ is a pivot: $p_{0}$ is assigned
the same position vector by both placements, which is what allows the
two response functions in \eqr{8.13} to be compared at a single point.
The pivot costs no generality, since translating $\mu$ leaves
$\bR=\Grad\vek{\lambda}$ untouched.}
\label{fig:ch8local}
\end{figure}

The pivot condition looks like an extra hypothesis and is not one. The
placement $\vek{\mu}$ is ours to choose, and replacing $\vek{\mu}$ by
$\vek{\mu}+\bc$ for a constant vector $\bc$ shifts every $\bY$ by $\bc$
and leaves $\bR=\Grad\vek{\lambda}$ completely unchanged, a constant
having zero gradient. So one may always translate $\mu$ until
$\vek{\lambda}(\bX_{0})=\bX_{0}$, and since $\bR$ is the only feature of
$\vek{\lambda}$ that will be used, nothing has been given up. What has
been bought is that both response functions may be evaluated at the same
label $\bX_{0}$, so that \eqr{8.13} compares like with like:
\begin{equation}
\begin{aligned}
   F(p_{0},t)
   &=F_{\kappa}(\bx,\bv,\theta,\bF_{\kappa},\bL,\grad\theta;\bX_{0})\\
   &=F_{\mu}(\bx,\bv,\theta,\bF_{\mu},\bL,\grad\theta;\bX_{0}),
   \qquad\text{with }\bF_{\kappa}=\bF_{\mu}\bR(\bX_{0}),
\end{aligned}
   \tag{8.13}\label{eq:8.13}
\end{equation}
all arguments being evaluated at $\bX_{0}$.

\subsection*{One experiment, two neighbourhoods}

Now the definition of material symmetry. Subject both
$N_{\kappa}(p_{0})$ and $N_{\mu}(p_{0})$ to the \emph{same} experiment,
consisting of a prescribed deformation $\bchi(\cdot,t)$ and temperature
distribution $\theta(\cdot,t)$. Acting on $N_{\kappa}(p_{0})$ the
experiment produces the fields
\begin{equation}
\begin{aligned}
   \bchi(\bX,t)&=\bchi(\bX_{0},t)+\bF(\bX_{0},t)(\bX-\bX_{0})
     +o(|\bX-\bX_{0}|),\\
   \hat\bv(\bX,t)&=\hat\bv(\bX_{0},t)+\dot\bF(\bX_{0},t)(\bX-\bX_{0})
     +o(|\bX-\bX_{0}|),\\
   \theta(\bX,t)&=\theta(\bX_{0},t)
     +\ip{\Grad\theta(\bX_{0},t)}{(\bX-\bX_{0})}+o(|\bX-\bX_{0}|),
\end{aligned}
   \tag{8.14}\label{eq:8.14}
\end{equation}
for $\bX\in N_{\kappa}(p_{0})$, with $\hat\bv(\bX_{0},t)=
\dot\bchi(\bX_{0},t)$ and $\bF(\bX_{0},t)=\Grad\bchi\rest{\bX_{0}}$;
and acting on $N_{\mu}(p_{0})$ it produces
\begin{equation}
\begin{aligned}
   \bchi(\bY,t)&=\bchi(\bX_{0},t)+\bF(\bX_{0},t)(\bY-\bX_{0})
     +o(|\bY-\bX_{0}|),\\
   \hat\bv(\bY,t)&=\hat\bv(\bX_{0},t)+\dot\bF(\bX_{0},t)(\bY-\bX_{0})
     +o(|\bY-\bX_{0}|),\\
   \theta(\bY,t)&=\theta(\bX_{0},t)
     +\ip{\Grad\theta(\bX_{0},t)}{(\bY-\bX_{0})}+o(|\bY-\bX_{0}|),
\end{aligned}
   \tag{8.15}\label{eq:8.15}
\end{equation}
for $\bY\in N_{\mu}(p_{0})$. The point of writing both out is that the
coefficients are literally the same symbols: the values of the fields
and of their gradients agree at $\bX_{0}$, which is possible only
because $\bX_{0}$ is a pivot and therefore belongs to both
neighbourhoods.

Now compute the response twice. Using $N_{\kappa}(p_{0})$ as reference,
it is $F_{\kappa}(\bx,\bv,\theta,\bF,\bL,\grad\theta;\bX_{0})$ with
\[
   \bx=\bchi(\bX_{0},t),\quad
   \bv=\dot\bchi(\bX_{0},t),\quad
   \theta=\theta(\bX_{0},t),\quad
   \bF=\bF(\bX_{0},t),
\]
\[
   \bL=\dot\bF(\bX_{0},t)\bF(\bX_{0},t)^{-1},\quad
   \grad\theta=\bF(\bX_{0},t)^{-t}\Grad\theta(\bX_{0},t),
\]
the last two being \eqr{4.3} and \eqr{8.6} read backwards. Using
$N_{\mu}(p_{0})$ as reference it is
$F_{\mu}(\bx,\bv,\theta,\bF,\bL,\grad\theta;\bX_{0})$, and by
\eqr{8.15} the six arguments are the very same six numbers, vectors and
tensors. So we have two functions evaluated at identical arguments.

Having subjected two different local neighbourhoods of $p_{0}$ to the
same experiment, we should expect the two response functions, evaluated
at the same arguments, to return \emph{different} values, and \eqr{8.13}
says exactly why: the two are related by an insertion of $\bR$ in the
fourth slot, and a function need not be blind to that. However, there
may be distinguished maps $\vek{\lambda}$ for which the two responses at
$p_{0}$ are nevertheless identical. For such a map no experiment
whatever can tell the two neighbourhoods apart. These are the
\emph{material symmetry transformations}.

Suppose then that $N_{\kappa}(p_{0})$ and $N_{\mu}(p_{0})$ are related
by a symmetry transformation, so that
\[
   F_{\kappa}(\bx,\bv,\theta,\bG,\bL,\grad\theta;\bX_{0})
   =F_{\mu}(\bx,\bv,\theta,\bG,\bL,\grad\theta;\bX_{0})
\]
for every admissible value $\bG$ of the fourth argument. Equation
\eqr{8.13} read as an identity between functions says
\[
   F_{\mu}(\bx,\bv,\theta,\bG,\bL,\grad\theta;\bX_{0})
   =F_{\kappa}(\bx,\bv,\theta,\bG\bR,\bL,\grad\theta;\bX_{0}),
\]
again for every $\bG$, since \eqr{8.13} holds whatever deformation the
material happens to be undergoing. Eliminating $F_{\mu}$ between the two
displays and renaming $\bG$ as $\bF$,
\begin{equation}
\begin{aligned}
   &F_{\kappa}(\bx,\bv,\theta,\bF,\bL,\grad\theta;\bX_{0})
   =F_{\kappa}(\bx,\bv,\theta,\bF\bR,\bL,\grad\theta;\bX_{0}),\\
   &\qquad\text{where }\bR=\Grad\vek{\lambda}(\bX)\rest{\bX_{0}} .
\end{aligned}
   \tag{8.16}\label{eq:8.16}
\end{equation}

Two features of \eqr{8.16} deserve emphasis before it is used.

It is a statement about \emph{one} function. The second reference
configuration has vanished from it entirely, which is what makes it
usable: \eqr{8.16} is a functional equation constraining $F_{\kappa}$,
and one may test a candidate constitutive function against it without
ever constructing $\mu$.

And it is required to hold \emph{for all admissible} $\bF$, with $\bR$
fixed. This quantifier is not decoration. Every argument in the rest of
the section, and the whole of Section~\ref{sec:ch8noll}, proceeds by
substituting some new tensor for $\bF$ in \eqr{8.16}, and each such
substitution is legal only because the equation was demanded of every
$\bF$ to begin with.

\begin{remark}[symmetry is a local notion]\label{rem:ch8localsym}
The map $\vek{\lambda}$ in this discussion is defined on a neighbourhood
of the material point in question, and on nothing more. In particular
$\bR$ need not be the gradient of a single function defined on the whole
body: there is no requirement that the local relabellings at different
points fit together into a global change of reference configuration.

Two things follow. Different points of one body may have different
symmetries, so the theory accommodates a composite, in which the type of
material varies from place to place, or a polycrystal, in which the
lattice orientation does. And, less obviously, the freedom to choose
$\vek{\lambda}$ point by point is exactly what Section~\ref{sec:ch8noll}
will exploit and what makes Noll's rule matter in plasticity, where the
local relaxed configurations at neighbouring points genuinely fail to
fit together.
\end{remark}

\subsection*{The symmetry group}

It follows from \eqr{8.16} that $\bF\bR$ is the value of an admissible
deformation gradient at $\bX_{0}$: the right-hand side evaluates the
constitutive function with $\bF\bR$ in the deformation-gradient slot,
and a constitutive function is only defined for arguments the material
can actually present. If $\kappa$ is an occupiable reference
configuration this requires
\[
   0<\det(\bF\bR)=(\det\bF)(\det\bR),
   \qquad\text{with }\det\bF>0,
\]
whence $\det\bR>0$. Denote by $\mathcal{G}_{\kappa(p_{0})}$ the set of
those $\bR$ for which \eqr{8.16} holds for all admissible values of
$\bF$:
\begin{equation}
   \mathcal{G}_{\kappa(p_{0})}
   =\{\bR:\det\bR>0\ \text{and \eqr{8.16} holds}\}.
   \tag{8.17}\label{eq:8.17}
\end{equation}

This set is a \emph{group}. Three properties have to be checked, and
each is a two-line calculation in which every step uses the ``for all
$\bF$'' clause of \eqr{8.16}. Throughout, all arguments not involving
$\bR$ are suppressed, so that $F_{\kappa}(\bF)$ abbreviates
$F_{\kappa}(\bx,\bv,\theta,\bF,\bL,\grad\theta;\bX_{0})$.

\medskip
\noindent
\textbf{(1) $\hI\in\mathcal{G}_{\kappa(p_{0})}$.} With $\bR=\hI$,
\eqr{8.16} reads $F_{\kappa}(\bF)=F_{\kappa}(\bF\hI)=F_{\kappa}(\bF)$,
which is true of any function whatever, and $\det\hI=1>0$. The identity
is a symmetry of every material: doing nothing to the reference
neighbourhood is certainly undetectable.

\medskip
\noindent
\textbf{(2) If $\bR_{1},\bR_{2}\in\mathcal{G}_{\kappa(p_{0})}$ then
$\bR_{1}\bR_{2}\in\mathcal{G}_{\kappa(p_{0})}$.} Compute, for arbitrary
admissible $\bF$,
\[
   F_{\kappa}\big(\bF(\bR_{1}\bR_{2})\big)
   =F_{\kappa}\big((\bF\bR_{1})\bR_{2}\big)
   =F_{\kappa}(\bF\bR_{1})
   =F_{\kappa}(\bF).
\]
The first equality is associativity of the tensor product, nothing more.
The second applies \eqr{8.16} for $\bR_{2}$ \emph{with $\bF$ replaced by
$\bF\bR_{1}$}, which is permitted because $\bF\bR_{1}$ is itself an
admissible deformation gradient and \eqr{8.16} was demanded for all such.
The third applies \eqr{8.16} for $\bR_{1}$. Finally
$\det(\bR_{1}\bR_{2})=(\det\bR_{1})(\det\bR_{2})>0$, so the product
qualifies on both counts.

\medskip
\noindent
\textbf{(3) If $\bR\in\mathcal{G}_{\kappa(p_{0})}$ then
$\bR^{-1}\in\mathcal{G}_{\kappa(p_{0})}$.} The inverse exists because
$\det\bR>0$. Apply \eqr{8.16} with $\bF$ replaced by $\bF\bR^{-1}$:
\[
   F_{\kappa}(\bF\bR^{-1})
   =F_{\kappa}\big((\bF\bR^{-1})\bR\big)
   =F_{\kappa}(\bF\hI)
   =F_{\kappa}(\bF),
\]
which is \eqr{8.16} for $\bR^{-1}$. And
$\det(\bR^{-1})=1/\det\bR>0$.

\medskip
Associativity needs no proof, being inherited from the composition of
linear maps, so $\mathcal{G}_{\kappa(p_{0})}$ is a subgroup of the group
of invertible tensors of positive determinant. The physical reading of
the three properties is worth pausing on: (1) says that doing nothing is
undetectable, (2) that two undetectable rearrangements performed in
succession are undetectable, and (3) that undoing an undetectable
rearrangement is undetectable. Stated that way the group structure is
obvious, which is the point; the calculations above are only the
translation of the obvious into the language of \eqr{8.16}.

\subsection*{The unimodular restriction}

Property (2) applied repeatedly gives $\bR^{n}\in
\mathcal{G}_{\kappa(p_{0})}$ for every positive integer $n$: if
$\bR^{n-1}$ is in the group then so is $\bR^{n-1}\bR=\bR^{n}$, and the
induction starts at $\bR^{1}=\bR$. Hence
$F_{\kappa}(\bF)=F_{\kappa}(\bF\bR^{n})$ for every $n$, and
\[
   \det(\bF\bR^{n})=(\det\bF)(\det\bR^{n})=(\det\bF)(\det\bR)^{n},
\]
using Theorem~\ref{thm:det} twice. Let $n\to\infty$ and read off the
two cases.

If $\det\bR>1$ then $(\det\bR)^{n}\to\infty$, so
$\det(\bF\bR^{n})\to\infty$. If $\det\bR<1$ then $(\det\bR)^{n}\to0$ and
$\det(\bF\bR^{n})\to0$. In the first case the material at $p_{0}$ is
dilated without bound and in the second compacted without bound, and in
neither case does its response change at all: the left-hand side
$F_{\kappa}(\bF)$ never moved.

That is untenable, and it is worth being explicit about why rather than
resting on the word. By \eqr{2.50} and \eqr{2.59} the current density is
$\rho=\rho_{\kappa}(\bX)/\det\bF$, so a deformation gradient with
$\det\bF\to0$ is one in which a finite mass is squeezed into vanishing
volume and $\rho\to\infty$, while $\det\bF\to\infty$ is unlimited
rarefaction. A material insensitive to either is a material one could
compress to a point at no cost in stress and no change in energy. No
substance behaves so. To avoid the conclusion we must restrict $\bR$ so
that neither limit can be reached, that is, so that $(\det\bR)^{n}$ does
not run away, and for a positive number that means $\det\bR=1$.
Therefore
\begin{equation}
   \mathcal{G}_{\kappa(p_{0})}\subset U=\{\bR:\det\bR=1\},
   \tag{8.18}\label{eq:8.18}
\end{equation}
which is called the \emph{unimodular group}. It is a group by the same
determinant identities used above: $\det\hI=1$, $\det(\bR_{1}\bR_{2})=
1\cdot1=1$ and $\det(\bR^{-1})=1/1=1$.

\begin{remark}[what \eqr{8.18} does and does not say]\label{rem:ch8unim}
It does not say that symmetry transformations preserve shape. It says
they preserve \emph{volume}, by Remark~\ref{rem:vol}, since the
determinant is the ratio of oriented volumes. A unimodular $\bR$ may
stretch the reference neighbourhood by a factor of ten in one direction
provided it compresses it correspondingly in others, and \eqr{8.18}
permits that; whether the material notices is a separate question,
answered by \eqr{8.19}.

Nor is $U$ the largest group one might have written. Some authors define
the unimodular group by $|\det\bR|=1$, which admits reflections. Here
$\det\bR>0$ was forced before \eqr{8.18} was reached, by the requirement
that $\bF\bR$ be an admissible deformation gradient, so reflections are
excluded from the outset. They are excluded for a physical reason and
not a conventional one: a reflection of the reference neighbourhood
would have to be realised by an actual motion of material, and no motion
can pass continuously from $\det\bF>0$ to $\det\bF<0$ without passing
through $\det\bF=0$, at which the deformation ceases to be invertible
(Remark~\ref{rem:localinv}).

Finally, \eqr{8.18} is an upper bound and nothing else. Every material
of the class considered satisfies it, in every reference configuration.
The interesting question is how far below the bound a given material
sits, and that is the question the next two definitions answer.
\end{remark}

\subsection*{Solid and fluid, defined by the symmetry group}

It was Noll's idea to classify materials as solid or fluid on the basis
of the symmetry group. The classification is not a description of how a
substance looks or pours; it is a statement about which rearrangements
of a material neighbourhood the substance can detect.

\medskip
\noindent\textbf{Fluids.} Take the simplest constitutive equation there
is: the stress in a compressible inviscid fluid under isothermal
conditions,
\[
   \bT=-p(\rho)\I ,
\]
where $p(\rho)$ is the constitutive function for the density-dependent
pressure. Ask which $\bR$ satisfy \eqr{8.16} for it. From \eqr{2.50} and
\eqr{2.59},
\[
   \rho=\frac{\rho_{\kappa}(\bX)}{\det\bF}
   =\frac{\rho_{\kappa}(\bX)}{\det(\bF\bR)}
   \qquad\text{provided }\bR\in U,
\]
the second equality holding because $\det(\bF\bR)=(\det\bF)(\det\bR)=
\det\bF$ when $\det\bR=1$. So replacing $\bF$ by $\bF\bR$ leaves $\rho$
unchanged, hence leaves $p(\rho)$ unchanged, hence leaves $\bT$
unchanged: \eqr{8.16} holds for every $\bR\in U$. Therefore
$U\subset\mathcal{G}_{\kappa(p_{0})}$, and with \eqr{8.18} in the
opposite direction,
\[
   \mathcal{G}_{\kappa(p_{0})}=U .
\]
This motivates the definition: a \emph{fluid} is a material of the class
considered here for which there exists a local reference configuration
$\kappa^{*}(p_{0})$ such that
$\mathcal{G}_{\kappa^{*}(p_{0})}=U$.

The stress depended on $\bF$ only through $\det\bF$, and $\det$ is
exactly the feature of $\bF$ that a unimodular map cannot change. A
fluid has the largest symmetry group available: it is a material that
cannot detect any volume-preserving rearrangement of its own reference
neighbourhood whatever. That is what it means to say a fluid has no
preferred shape, expressed constitutively rather than anecdotally.

\medskip
\noindent\textbf{Solids.} For these the situation is different, and the
difference is one sentence of physics: experience indicates that a
strain induced by the map $\vek{\lambda}$ is experimentally detectable.
A solid remembers a shape, and deforming it away from that shape
produces stress that can be measured.

If $\vek{\lambda}$ is to be a symmetry transformation, and hence
undetectable, it must therefore be strain free. Strain was defined in
Chapter 2 through the right Cauchy--Green tensor, \eqr{2.96} and
\eqr{2.97}: a deformation with gradient $\bR$ is strain free exactly
when the associated $\bC=\bR^{t}\bR$ equals $\hI$, for then the Lagrange
strain $\bE=\tfrac12(\bC-\hI)$ vanishes and, by
Proposition~\ref{prop:undeformed}, every material line element keeps its
length and every pair of them its angle. So a symmetry transformation of
a solid satisfies $\bR^{t}\bR=\hI$. Combined with \eqr{8.18} this makes
$\bR$ a rotation: $\bR^{t}\bR=\hI$ makes $\bR$ orthogonal, hence
$\det\bR=\pm1$ by Proposition~\ref{prop:orth}, and $\det\bR=1$ selects
the proper orthogonal case.

Accordingly a \emph{solid} is defined to be a material of the class
considered for which there exists a local reference $\kappa^{*}(p_{0})$
such that $\mathcal{G}_{\kappa^{*}(p_{0})}$ is contained in
$\Orth^{+}$, the set of proper-orthogonal tensors, or simply the
rotations:
\begin{equation}
   \mathcal{G}_{\kappa^{*}(p_{0})}\subset\Orth^{+}
   =\{\bR:\bR\in U\ \text{and}\ \bR^{t}\bR=\hI\}.
   \tag{8.19}\label{eq:8.19}
\end{equation}

We note that $\kappa^{*}(p_{0})$ need bear no relationship to any
particular reference configuration one may use for the purposes of
analysis. In a solid it would correspond to a unit cell of an undistorted
lattice of a crystalline material. Such cells are mapped to themselves
by certain discrete rotations, and those rotations are what characterize
the crystalline symmetry at hand: four of them about a four-fold axis
for a square cell, twenty-four in all for a cubic crystal, and only the
identity for a triclinic one. For this reason $\kappa^{*}(p_{0})$ is
often called an \emph{undistorted} local configuration of the solid.

\begin{remark}[isotropy, and the whole ladder in one line]
\label{rem:ch8ladder}
Definitions \eqr{8.18} and \eqr{8.19} between them organise every
material of the class considered on a single scale, and the scale is
worth writing out because the later chapters use its rungs by name.
Every symmetry group satisfies
\[
   \{\hI\}\ \subseteq\ \mathcal{G}_{\kappa(p_{0})}\ \subseteq\ U ,
\]
with the lower end forced by property (1) above and the upper by
\eqr{8.18}. Between the ends:
\begin{itemize}\itemsep1pt
\item $\mathcal{G}=\{\hI\}$: no symmetry at all. A triclinic crystal.
      Every rearrangement, however slight, is detectable.
\item $\mathcal{G}$ a finite subgroup of $\Orth^{+}$: a crystalline
      solid, the finite group being its point group.
\item $\mathcal{G}\subseteq\Orth^{+}$, possibly infinite: a solid.
      A transversely isotropic solid, for instance, has the continuous
      group of rotations about one axis.
\item $\mathcal{G}=\Orth^{+}$: an \emph{isotropic} solid. It is the
      largest a solid's group can be, since \eqr{8.19} caps it, so an
      isotropic solid is a solid with as much symmetry as the definition
      of a solid allows: no direction in $\kappa^{*}$ is distinguished
      from any other.
\item $\mathcal{G}=U$: a fluid.
\end{itemize}
Two consequences follow immediately from $\Orth^{+}\subset U$, and both
are the kind of statement that is obvious once the group language is in
place and obscure without it.

\emph{Every fluid is isotropic.} A fluid's group is $U$, which contains
$\Orth^{+}$; so a fluid has at least the symmetry an isotropic solid
has. The converse fails: an isotropic solid has $\Orth^{+}$ and not $U$,
and the elements of $U$ outside $\Orth^{+}$ are exactly the
volume-preserving distortions it can feel and a fluid cannot. Steel and
water are both indifferent to direction; only one of them is indifferent
to shape.

\emph{The classification is exhaustive but not complete.} A group
strictly between $\Orth^{+}$ and $U$ describes a material that is
neither a solid nor a fluid in the above sense. Such groups exist and
such materials exist, liquid crystals among them. The two definitions
are the two extremes of a range, and the range in between is where much
of modern constitutive theory lives.
\end{remark}

\begin{figure}[htbp]\centering
\renewcommand{\thefigure}{A}
\begin{tikzpicture}[scale=1.0,
   ar/.style={-{Stealth[length=2.3mm]},thick},
   dot/.style={fill=black}]
% =============== ROW 1 : a solid ===============
\node[font=\small,anchor=west] at (-7.40,1.30) {\textsc{a solid}};
% lattice 1 : square, undistorted
\foreach \i in {-1,0,1}{\foreach \j in {-1,0,1}{
   \fill (-5.60+0.62*\i,0.62*\j) circle (1.3pt);}}
\draw[thin,gray!70] (-6.22,-0.62) rectangle (-5.60,0);
\node[font=\footnotesize,align=center,anchor=north] at (-5.60,-1.15)
   {$\kappa^{*}(p_{0})$\\ undistorted};
% arrow 1
\draw[ar] (-4.60,0) -- (-2.30,0);
\node[font=\footnotesize,anchor=south] at (-3.45,0.14)
   {$\bR\in\Orth^{+}$};
\node[font=\footnotesize,anchor=north,text=gray!80] at (-3.45,-0.16)
   {$90^{\circ}$ turn};
% lattice 2 : square again
\foreach \i in {-1,0,1}{\foreach \j in {-1,0,1}{
   \fill (-1.30+0.62*\i,0.62*\j) circle (1.3pt);}}
\draw[thin,gray!70] (-1.92,-0.62) rectangle (-1.30,0);
\node[font=\footnotesize,align=center,anchor=north] at (-1.30,-1.15)
   {the same lattice\\ undetectable};
% arrow 2
\draw[ar] (-0.30,0) -- (2.00,0);
\node[font=\footnotesize,anchor=south] at (0.85,0.14) {$\bH\in U$};
\node[font=\footnotesize,anchor=north,text=gray!80] at (0.85,-0.16)
   {$\bH^{t}\bH\neq\hI$};
% lattice 3 : rectangular
\foreach \i in {-1,0,1}{\foreach \j in {-1,0,1}{
   \fill (3.35+0.99*\i,0.39*\j) circle (1.3pt);}}
\draw[thin,gray!70] (2.36,-0.39) rectangle (3.35,0);
\node[font=\footnotesize,align=center,anchor=north] at (3.35,-1.15)
   {a strained lattice\\ detectable};
% =============== ROW 2 : a fluid ===============
\node[font=\small,anchor=west] at (-7.40,-2.50) {\textsc{a fluid}};
\draw[thick,fill=blue!7] (-6.25,-4.25) rectangle (-4.95,-2.95);
\foreach \p in {(-6.02,-4.03),(-5.44,-3.72),(-5.90,-3.28),
                (-5.18,-4.08),(-5.52,-3.14),(-5.08,-3.50)}{
   \fill[gray!75] \p circle (1.1pt);}
\node[font=\footnotesize,anchor=north] at (-5.60,-4.45) {volume $V$};
\draw[ar] (-4.60,-3.60) -- (-2.30,-3.60);
\node[font=\footnotesize,anchor=south] at (-3.45,-3.46)
   {any $\bH\in U$};
\draw[thick,fill=blue!7] (-1.95,-4.25)--(-0.25,-4.25)--(0.30,-3.26)
   --(-1.40,-3.26)--cycle;
\foreach \p in {(-1.62,-4.03),(-0.92,-3.83),(-1.20,-3.52),
                (-0.55,-4.11),(-0.48,-3.44),(-0.10,-3.72)}{
   \fill[gray!75] \p circle (1.1pt);}
\node[font=\footnotesize,anchor=north] at (-0.85,-4.45) {volume $V$};
\node[font=\footnotesize,align=left,anchor=west] at (0.95,-3.60)
   {same $V\Rightarrow$ same $\rho$\\
    $\Rightarrow$ same $\bT=-p(\rho)\I$\\
    undetectable, so $\mathcal{G}=U$};
% =============== ROW 3 : the ladder ===============
\draw[thick,rounded corners=5pt] (-6.60,-7.70) rectangle (1.00,-5.75);
\node[font=\footnotesize,anchor=north west] at (-6.45,-5.85) {$U$};
\draw[thick,rounded corners=5pt] (-4.70,-7.40) rectangle (-0.30,-6.20);
\node[font=\footnotesize,anchor=north west] at (-4.55,-6.28)
   {$\Orth^{+}$};
\draw[thick,rounded corners=4pt] (-2.60,-7.20) rectangle (-0.95,-6.65);
\node[font=\footnotesize] at (-1.78,-6.92) {$\mathcal{G}$};
\node[font=\footnotesize,anchor=west,text=gray!85] at (1.90,-6.05)
   {fluid: $\mathcal{G}=U$};
\node[font=\footnotesize,anchor=west,text=gray!85] at (1.90,-6.72)
   {isotropic solid: $\mathcal{G}=\Orth^{+}$};
\node[font=\footnotesize,anchor=west,text=gray!85] at (1.90,-7.40)
   {crystal: $\mathcal{G}$ finite};
\draw[gray!60,-{Stealth[length=1.6mm]}] (1.82,-6.05) -- (1.08,-5.95);
\draw[gray!60,-{Stealth[length=1.6mm]}] (1.82,-6.72) -- (-0.22,-6.60);
\draw[gray!60,-{Stealth[length=1.6mm]}] (1.82,-7.40) -- (-0.88,-7.05);
\end{tikzpicture}
\caption{Solid and fluid, told apart by the symmetry group. Top: in an
undistorted configuration a crystalline solid is carried onto itself by
certain rotations, which no experiment can detect, but a unimodular map
that is not a rotation strains the lattice, and by \eqr{8.19} that is
detectable; the group is confined to $\Orth^{+}$. Middle: the response
of an inviscid fluid depends on $\bF$ only through $\det\bF$, and no
member of $U$ changes that, so every volume-preserving rearrangement is
undetectable and the group is the whole of $U$. Bottom: the two
definitions as nested sets, with $\Orth^{+}\subset U$ showing at a
glance that every fluid is isotropic while an isotropic solid is not a
fluid.}
\label{fig:ch8solidfluid}
\end{figure}

%=======================================================================
\section{Noll's rule}\label{sec:ch8noll}
%=======================================================================

Material symmetry, as \eqr{8.16} defines it, is a characteristic of two
things at once: of the material, and of the reference configuration
adopted. The definitions of solid and fluid both had to say ``there
exists a local reference configuration such that\ldots'', which is an
admission that the group could look different in another one. It is
therefore of interest to determine exactly how much the choice of
reference matters. The answer is complete and is due to Noll.

Consider two reference configurations $\kappa_{1}$ and $\kappa_{2}$ in
which the positions of $p_{0}\in B$ are $\bX_{1}$ and $\bX_{2}$
respectively, and assume a smooth invertible map $\vek{\pi}$, defined on
a neighbourhood of $\bX_{1}$, with $\bX_{2}=\vek{\pi}(\bX_{1})$. Write
$\bK=\Grad\vek{\pi}\rest{\bX_{1}}$. This is the situation of \eqr{8.8}
with new names, so \eqr{8.12} applies verbatim; suppressing the
variables that do not involve a reference configuration, it reads
\begin{equation}
\begin{aligned}
   F_{\kappa_{2}}(\bF;\bX_{2})&=F_{\kappa_{1}}(\bF\bK;\bX_{1}),
   \qquad\text{where }\bK\\
   &=\Grad\vek{\pi}\rest{\bX_{1}} .
\end{aligned}
   \tag{8.20}\label{eq:8.20}
\end{equation}
Here $\bF$ in the argument slot is the deformation gradient computed
from $\kappa_{2}$, and $\bF\bK$ is the corresponding one computed from
$\kappa_{1}$, which is \eqr{8.9} again.

Everything that follows is obtained by putting different tensors into
the one slot of \eqr{8.20}, so it is worth repeating that this is legal:
\eqr{8.20} holds for every admissible $\bF$, and $\det\bK>0$ because
both configurations are occupiable, so $\bF\bK$, $\bF\bK^{-1}$ and
$\bF\bR\bK^{-1}$ are all admissible whenever $\bF$ is.

Now replace $\bF$ by $\bF\bK^{-1}$ in \eqr{8.20}. The right-hand side
becomes $F_{\kappa_{1}}(\bF\bK^{-1}\bK;\bX_{1})=F_{\kappa_{1}}
(\bF\hI;\bX_{1})=F_{\kappa_{1}}(\bF;\bX_{1})$, so
\begin{equation}
   F_{\kappa_{2}}(\bF\bK^{-1};\bX_{2})=F_{\kappa_{1}}(\bF;\bX_{1}).
   \tag{8.21}\label{eq:8.21}
\end{equation}
This is \eqr{8.20} solved the other way round, and it is the form we
shall use, because it has $\kappa_{1}$, whose symmetry group is known,
standing alone on the right.

Now suppose $\bR\in\mathcal{G}_{\kappa_{1}(p_{0})}$, so that by
\eqr{8.16}
\begin{equation}
   F_{\kappa_{1}}(\bF;\bX_{1})=F_{\kappa_{1}}(\bF\bR;\bX_{1})
   \tag{8.22}\label{eq:8.22}
\end{equation}
for every admissible $\bF$. Replacing $\bF$ by $\bF\bR$ in \eqr{8.21},
\begin{equation}
   F_{\kappa_{1}}(\bF\bR;\bX_{1})
   =F_{\kappa_{2}}(\bF\bR\bK^{-1};\bX_{2}).
   \tag{8.23}\label{eq:8.23}
\end{equation}
Chaining \eqr{8.23}, \eqr{8.22} and \eqr{8.21} in that order gives
\begin{equation}
   F_{\kappa_{2}}(\bF\bR\bK^{-1};\bX_{2})
   =F_{\kappa_{1}}(\bF;\bX_{1})
   =F_{\kappa_{2}}(\bF\bK^{-1};\bX_{2}).
   \tag{8.24}\label{eq:8.24}
\end{equation}
The two ends of \eqr{8.24} are values of the same function
$F_{\kappa_{2}}$, so the statement is now entirely about $\kappa_{2}$:
$\kappa_{1}$ has served its purpose and may be discarded. All that
remains is to recognise \eqr{8.16} in it.

Set $\hat\bF=\bF\bK^{-1}$, so that $\bF=\hat\bF\bK$ and
\[
   \bF\bR\bK^{-1}=\hat\bF\bK\bR\bK^{-1}.
\]
Then \eqr{8.24} becomes
\begin{equation}
   F_{\kappa_{2}}(\hat\bF;\bX_{2})
   =F_{\kappa_{2}}(\hat\bF\bK\bR\bK^{-1};\bX_{2}).
   \tag{8.25}\label{eq:8.25}
\end{equation}
As $\bF$ runs over all admissible tensors so does $\hat\bF=\bF\bK^{-1}$,
because $\bF\mapsto\bF\bK^{-1}$ is a bijection of the admissible set
onto itself; so \eqr{8.25} holds for every admissible $\hat\bF$. Comparing
it with \eqr{8.16}, and noting
\[
   \det(\bK\bR\bK^{-1})
   =(\det\bK)(\det\bR)(\det\bK)^{-1}=\det\bR>0,
\]
we conclude that $\bK\bR\bK^{-1}\in\mathcal{G}_{\kappa_{2}(p_{0})}$. We
thus arrive at \emph{Noll's rule}:
\begin{equation}
   \mathcal{G}_{\kappa_{2}(p_{0})}
   =\{\bK\bR\bK^{-1}:\bR\in\mathcal{G}_{\kappa_{1}(p_{0})}\}.
   \tag{8.26}\label{eq:8.26}
\end{equation}

\begin{remark}[the missing half of \eqr{8.26}]\label{rem:ch8converse}
What the calculation established is one inclusion,
$\bK\mathcal{G}_{\kappa_{1}(p_{0})}\bK^{-1}\subseteq
\mathcal{G}_{\kappa_{2}(p_{0})}$: every conjugate of a $\kappa_{1}$
symmetry is a $\kappa_{2}$ symmetry. Equality needs the other direction,
and it costs one line.

The situation is symmetric in the two configurations. Running the whole
argument again with $\kappa_{1}$ and $\kappa_{2}$ interchanged uses
$\vek{\pi}^{-1}$ in place of $\vek{\pi}$, whose gradient at $\bX_{2}$ is
$\bK^{-1}$, and yields
$\bK^{-1}\mathcal{G}_{\kappa_{2}(p_{0})}\bK\subseteq
\mathcal{G}_{\kappa_{1}(p_{0})}$. So if $\bS\in
\mathcal{G}_{\kappa_{2}(p_{0})}$ then $\bK^{-1}\bS\bK\in
\mathcal{G}_{\kappa_{1}(p_{0})}$, and
\[
   \bS=\bK\big(\bK^{-1}\bS\bK\big)\bK^{-1}
\]
exhibits $\bS$ as a conjugate of a member of
$\mathcal{G}_{\kappa_{1}(p_{0})}$. Both inclusions hold, so \eqr{8.26}
is an equality. The gap is worth closing because the two statements are
not interchangeable: the inclusion alone would leave open the
possibility that $\kappa_{2}$ has \emph{more} symmetry than $\kappa_{1}$,
and the whole force of the section is that it has exactly as much.
\end{remark}

\subsection*{Why conjugation, and why nothing else would do}

Equation \eqr{8.26} is stated compactly and can be read as a formula
that fell out of an algebraic manipulation. It is not. Conjugation is
the only thing that could have appeared, and there are three ways to see
that, each of which is worth having.

\medskip
\noindent\textbf{First: a symmetry is an object, and $\bK$ is a change
of description.} Recall what $\bR$ is. By \eqr{8.16} and the pivot
arrangement of Figure~\ref{fig:ch8local}, $\bR=\Grad\vek{\lambda}$ where
$\vek{\lambda}$ carries a neighbourhood of $\bX_{0}$ to another
neighbourhood of the same point. So $\bR$ is a linear map of the local
reference translation space to \emph{itself}: it takes a material fibre
of the reference neighbourhood and returns another material fibre of the
same neighbourhood, and it is a symmetry when the material cannot tell
the two apart.

Now change the reference. The tensor $\bK$ maps $\hVt_{\kappa_{1}}$ to
$\hVt_{\kappa_{2}}$: it is the dictionary between the two descriptions
of the same set of material fibres. A fibre described in $\kappa_{2}$ by
$\dl\bX_{2}$ is described in $\kappa_{1}$ by
$\dl\bX_{1}=\bK^{-1}\dl\bX_{2}$. To apply the rearrangement, translate
into $\kappa_{1}$, apply $\bR$ there, and translate back:
\[
   \dl\bX_{2}\ \overset{\bK^{-1}}{\longmapsto}\ \bK^{-1}\dl\bX_{2}
   \ \overset{\bR}{\longmapsto}\ \bR\bK^{-1}\dl\bX_{2}
   \ \overset{\bK}{\longmapsto}\ \bK\bR\bK^{-1}\dl\bX_{2},
\]
which is Figure~\ref{fig:ch8noll}. That is the whole of \eqr{8.26}, and
it is the same statement as Theorem~\ref{thm:tenstrans}, the
transformation law for the components of a tensor under a change of
basis, with one difference: there the change of basis was orthogonal and
the conjugation read $\bQ\bA\bQ^{t}$; here $\bK$ is merely invertible
and the conjugation reads $\bK\bA\bK^{-1}$. Remark~\ref{rem:2ptlaw} made
the point that the transformation law identifies the type of an object.
Equation \eqr{8.26} is that law, and it identifies $\bR$ as a linear map
of a translation space to itself.

\medskip
\noindent\textbf{Second: any other rule would break under composition.}
Change the reference twice, from $\kappa_{1}$ to $\kappa_{2}$ with
gradient $\bK_{12}$, then from $\kappa_{2}$ to $\kappa_{3}$ with
gradient $\bK_{23}$. The composite map has gradient
$\bK_{13}=\bK_{23}\bK_{12}$, by the chain rule. Applying \eqr{8.26}
twice,
\[
   \bK_{23}\big(\bK_{12}\,\mathcal{G}_{\kappa_{1}(p_{0})}\,
   \bK_{12}^{-1}\big)\bK_{23}^{-1}
   =(\bK_{23}\bK_{12})\,\mathcal{G}_{\kappa_{1}(p_{0})}\,
   (\bK_{23}\bK_{12})^{-1}
   =\bK_{13}\,\mathcal{G}_{\kappa_{1}(p_{0})}\,\bK_{13}^{-1},
\]
using $(\bK_{23}\bK_{12})^{-1}=\bK_{12}^{-1}\bK_{23}^{-1}$, which is the
same as applying \eqr{8.26} once with $\bK_{13}$. Doing it in two steps
agrees with doing it in one. A rule such as
$\bR\mapsto\bK\bR\bK^{t}$ or $\bR\mapsto\bK^{t}\bR\bK$ fails this test
at once, since $\bK^{t}$ does not compose in the required way, and it
would then matter through which intermediate configurations one had
passed. Nothing physical can depend on that.

\medskip
\noindent\textbf{Third: conjugation is an isomorphism, and that is the
real content.} The map $\bR\mapsto\bK\bR\bK^{-1}$ is one to one, is onto
$\mathcal{G}_{\kappa_{2}(p_{0})}$ by \eqr{8.26}, and respects products:
\[
   (\bK\bR_{1}\bK^{-1})(\bK\bR_{2}\bK^{-1})
   =\bK\bR_{1}(\bK^{-1}\bK)\bR_{2}\bK^{-1}
   =\bK(\bR_{1}\bR_{2})\bK^{-1}.
\]
So the two symmetry groups are isomorphic as groups. Whatever is
intrinsic to the group survives a change of reference: how many elements
it has, whether it is finite, the order of each element, whether it is
abelian. What does \emph{not} survive is how the group sits inside the
set of all tensors, which is to say whether its elements happen to be
rotations. The distinction between the abstract group and its
representation is exactly the distinction between what a material is and
how a particular reference configuration displays it.

\begin{figure}[htbp]\centering
\renewcommand{\thefigure}{B}
\begin{tikzpicture}[scale=1.0,
   ar/.style={-{Stealth[length=2.4mm]},thick},
   nd/.style={inner sep=2pt}]
\node[nd] (a1) at (0,2.20)   {$\hVt_{\kappa_{2}}$};
\node[nd] (b1) at (5.40,2.20) {$\hVt_{\kappa_{2}}$};
\node[nd] (a2) at (0,0)      {$\hVt_{\kappa_{1}}$};
\node[nd] (b2) at (5.40,0)   {$\hVt_{\kappa_{1}}$};
\draw[ar] (a1)--(a2) node[midway,left,font=\small] {$\bK^{-1}$};
\draw[ar] (a2)--(b2) node[midway,below,font=\small] {$\bR$};
\draw[ar] (b2)--(b1) node[midway,right,font=\small] {$\bK$};
\draw[ar] (a1)--(b1)
   node[midway,above,font=\small] {$\bK\bR\bK^{-1}$};
\node[font=\small,text=gray!75] at (2.70,1.10) {commutes};
\node[font=\footnotesize,text=gray!85,align=center,anchor=north]
   at (2.70,-0.85)
   {to act in $\kappa_{2}$: go back to $\kappa_{1}$,\\
    act there, come forward again};
\end{tikzpicture}
\caption{Noll's rule as a change of description. A symmetry
transformation is a linear map of the local reference translation space
to itself, and $\bK$ is the dictionary between the two descriptions of
that space. Going round the square the long way and the short way must
give the same answer, and that single requirement is \eqr{8.26}. Since
$\det(\bK\bR\bK^{-1})=\det\bR$, the top arrow is unimodular exactly when
the bottom one is; but $\bK\bR\bK^{-1}$ is orthogonal only under special
circumstances, which is why $U$ is preserved and $\Orth^{+}$ is not.}
\label{fig:ch8noll}
\end{figure}

\begin{remark}[$\bK$ is local, and what that buys]\label{rem:ch8Klocal}
As with $\bR$ in Remark~\ref{rem:ch8localsym}, $\bK$ is the gradient at
$\bX_{1}$ of a map defined on a neighbourhood of $\bX_{1}$. It need not
be the gradient of any map defined on the whole body, and in the
applications that matter most it is not.

Plasticity is the standard example. There one imagines the deformation
gradient factored as $\bF=\bF_{e}\bF_{p}$, where $\bF_{p}$ carries the
reference neighbourhood to a locally relaxed state, from which $\bF_{e}$
produces the actual current state elastically. The relaxed states are
precisely local reference configurations in the sense of this section,
one for each material point, chosen so that the lattice is undistorted
there. In general no single configuration of the body relaxes every
neighbourhood at once: the relaxed neighbourhoods do not fit together,
$\bF_{p}$ is not a gradient, and the residual stresses of a plastically
worked body are the visible consequence. Noll's rule is what governs how
the symmetry group is carried along by $\bF_{p}$ as the material yields,
and it is the reason the theory can speak of a lattice orientation that
evolves while the crystal class does not.
\end{remark}

\subsection*{A worked example: a four-fold axis in a distorted
configuration}

Take the smallest interesting symmetry group and watch what \eqr{8.26}
does to it. Let $\{\hbE_{A}\}$ be a fixed orthonormal basis of $\hVt$
and let $\kappa_{1}$ be a configuration in which the material is a
lattice generated by $\hbE_{1}$ and $\hbE_{2}$, with a four-fold axis
along $\hbE_{3}$. Its symmetry group is cyclic of order four,
\[
   \mathcal{G}_{\kappa_{1}(p_{0})}
   =\{\hI,\bR,\bR^{2},\bR^{3}\},
   \qquad
   \bR\hbE_{1}=\hbE_{2},\quad
   \bR\hbE_{2}=-\hbE_{1},\quad
   \bR\hbE_{3}=\hbE_{3},
\]
so that $\bR$ is the rotation through $90^{\circ}$ about $\hbE_{3}$,
with components
\[
   [\bR]=\begin{pmatrix}0&-1&0\\ 1&0&0\\ 0&0&1\end{pmatrix}.
\]
Since $\bR^{t}\bR=\hI$ and $\det\bR=1$, this group lies in $\Orth^{+}$
and $\kappa_{1}$ is an undistorted configuration in the sense of
\eqr{8.19}: the material is a solid and $\kappa_{1}$ displays it as one.

Now change the reference by a pure stretch that preserves volume,
\[
   \bK=\alpha\,\hbE_{1}\otimes\hbE_{1}
      +\alpha^{-1}\hbE_{2}\otimes\hbE_{2}
      +\hbE_{3}\otimes\hbE_{3},
   \qquad
   [\bK]=\begin{pmatrix}\alpha&0&0\\ 0&\alpha^{-1}&0\\
                        0&0&1\end{pmatrix},
\]
with $\alpha>0$ and $\alpha\neq1$; note $\det\bK=1$, so $\kappa_{2}$
holds the same volume of material as $\kappa_{1}$. Its inverse is
$[\bK^{-1}]=\operatorname{diag}(\alpha^{-1},\alpha,1)$. Conjugate,
multiplying the matrices in order:
\[
   [\bK][\bR]
   =\begin{pmatrix}\alpha&0&0\\ 0&\alpha^{-1}&0\\ 0&0&1\end{pmatrix}
    \begin{pmatrix}0&-1&0\\ 1&0&0\\ 0&0&1\end{pmatrix}
   =\begin{pmatrix}0&-\alpha&0\\ \alpha^{-1}&0&0\\
                   0&0&1\end{pmatrix},
\]
and then
\[
   [\bH]:=[\bK][\bR][\bK^{-1}]
   =\begin{pmatrix}0&-\alpha&0\\ \alpha^{-1}&0&0\\
                   0&0&1\end{pmatrix}
    \begin{pmatrix}\alpha^{-1}&0&0\\ 0&\alpha&0\\
                   0&0&1\end{pmatrix}
   =\begin{pmatrix}0&-\alpha^{2}&0\\ \alpha^{-2}&0&0\\
                   0&0&1\end{pmatrix}.
\]
Now read off what has and has not changed.

\emph{The determinant is unchanged.} $\det\bH=
-(-\alpha^{2})(\alpha^{-2})\cdot 1$ computed along the third row gives
$\det\bH=1$, so $\bH\in U$, as \eqr{8.26} guarantees in general.

\emph{Orthogonality is destroyed.}
\[
   [\bH^{t}][\bH]
   =\begin{pmatrix}0&\alpha^{-2}&0\\ -\alpha^{2}&0&0\\
                   0&0&1\end{pmatrix}
    \begin{pmatrix}0&-\alpha^{2}&0\\ \alpha^{-2}&0&0\\
                   0&0&1\end{pmatrix}
   =\begin{pmatrix}\alpha^{-4}&0&0\\ 0&\alpha^{4}&0\\
                   0&0&1\end{pmatrix},
\]
which is $\hI$ only when $\alpha=1$. So
$\mathcal{G}_{\kappa_{2}(p_{0})}\not\subset\Orth^{+}$: in $\kappa_{2}$
the symmetry is no longer a rotation, and $\kappa_{2}$ is a distorted
configuration.

\emph{The group structure is untouched.} Squaring,
\[
   [\bH]^{2}
   =\begin{pmatrix}0&-\alpha^{2}&0\\ \alpha^{-2}&0&0\\
                   0&0&1\end{pmatrix}^{2}
   =\begin{pmatrix}-1&0&0\\ 0&-1&0\\ 0&0&1\end{pmatrix}
   =[\bR^{2}],
\]
the two factors of $\alpha$ cancelling, and $\bH^{4}=(\bH^{2})^{2}=\hI$.
So $\bH$ still has order four and
$\mathcal{G}_{\kappa_{2}(p_{0})}=\{\hI,\bH,\bH^{2},\bH^{3}\}$ is still
cyclic of order four, as Remark~\ref{rem:ch8converse} and the
isomorphism argument require. The characteristic polynomial is likewise
preserved by conjugation, so $\bH$ has the same eigenvalues $1$ and
$\pm i$ as the rotation it came from, even though it is not a rotation.

\emph{And the geometry is exactly what one would draw.} In $\kappa_{2}$
the lattice is generated by $\bK\hbE_{1}=\alpha\hbE_{1}$ and
$\bK\hbE_{2}=\alpha^{-1}\hbE_{2}$, a rectangle rather than a square.
Apply $\bH$ to those generators:
\[
   \bH(\alpha\hbE_{1})=\alpha\cdot\alpha^{-2}\hbE_{2}
   =\alpha^{-1}\hbE_{2},
   \qquad
   \bH(\alpha^{-1}\hbE_{2})=\alpha^{-1}\cdot(-\alpha^{2})\hbE_{1}
   =-\alpha\hbE_{1}.
\]
Each generator is carried onto the other, so $\bH$ maps the rectangular
lattice onto itself, exactly as $\bR$ mapped the square one onto itself.
The symmetry is still there and is still detectable by no experiment;
what has changed is that expressing it now requires a map that stretches
one lattice direction by $\alpha^{-2}$ while compressing the other by
$\alpha^{2}$. A rotation could not do it, because the two lattice
directions no longer have equal spacing.

\subsection*{Consequences: fluidity is intrinsic, undistortedness is
not}

Two conclusions follow from \eqr{8.26}, and they are the reason the
definitions of Section~\ref{sec:ch8symm} were phrased as they were.

\medskip
\noindent\textbf{$\Orth^{+}$ is not preserved.} The worked example is a
counterexample: $\mathcal{G}_{\kappa_{1}(p_{0})}\subset\Orth^{+}$ while
$\mathcal{G}_{\kappa_{2}(p_{0})}$ is not. So if $\kappa_{1}(p_{0})$ is
an undistorted state of a crystalline lattice, the lattice is distorted
in $\kappa_{2}(p_{0})$, and \eqr{8.19} fails there. This is why the
existence of \emph{some} local configuration in which \eqr{8.19} holds
is an essential part of the definition of a solid: solidity is the
statement that an undistorted configuration exists, not that the one in
front of us is undistorted.

It is worth recording exactly when nothing goes wrong. If
$\bK=c\bQ$ with $c>0$ and $\bQ$ orthogonal, then $\bK^{-1}=c^{-1}\bQ^{t}$
and
\[
   \bK\bR\bK^{-1}=c\bQ\bR\,c^{-1}\bQ^{t}=\bQ\bR\bQ^{t},
\]
the scalars cancelling, and $\bQ\bR\bQ^{t}$ is orthogonal whenever $\bR$
is, since $(\bQ\bR\bQ^{t})^{t}(\bQ\bR\bQ^{t})=
\bQ\bR^{t}\bQ^{t}\bQ\bR\bQ^{t}=\bQ\bR^{t}\bR\bQ^{t}=\bQ\bQ^{t}=\hI$.
So rotating an undistorted configuration, or dilating it uniformly,
leaves it undistorted, and the undistorted configurations of a solid
form a family closed under those operations and no others. Anything that
strains the neighbourhood spoils it, which is the same statement as
\eqr{8.19} viewed from the other side.

\medskip
\noindent\textbf{$U$ is preserved.} Suppose
$\mathcal{G}_{\kappa_{1}(p_{0})}=U$. Every $\bR\in U$ has
$\det(\bK\bR\bK^{-1})=\det\bR=1$, so $\bK U\bK^{-1}\subseteq U$;
conversely every $\bS\in U$ is $\bK(\bK^{-1}\bS\bK)\bK^{-1}$ with
$\det(\bK^{-1}\bS\bK)=\det\bS=1$, so $U\subseteq\bK U\bK^{-1}$. Hence
$\bK U\bK^{-1}=U$ and, by \eqr{8.26},
$\mathcal{G}_{\kappa_{2}(p_{0})}=U$ as well.

So \eqr{8.18} is satisfied in all configurations, and more: the
\emph{identification} of the symmetry group with the unimodular group is
independent of the reference configuration. Being a fluid is a property
of the substance and of nothing else, and no change of reference,
however violent, can make a fluid look like a solid or a solid like a
fluid. That is in accord with the intuitive notion of fluidity, and it
is the reason the definition of a fluid could have been stated without
the words ``there exists a local reference configuration'', whereas the
definition of a solid could not.

%=======================================================================
\section{Frame invariance}\label{sec:ch8frame}
%=======================================================================

In building the class of constitutive equations of
Section~\ref{sec:ch8elastic} we adopted the perspective of one
particular observer $O$. Chapter 7 was written to make the second
observer available, and this section spends it.

Recall the setting. An observer is equipped with a point space and its
translation space, and two observers agree about time intervals and
about distances between material points. Section 7.1 showed that these
two agreements force the relation between their perceptions to be
\eqr{7.3},
\[
   \bx^{+}=\bQ(t)\bx+\bc(t),
\]
with $\bc(t)$ arbitrary and $\bQ(t)$ a spatially uniform, time-dependent
two-point orthogonal tensor, together with $t^{+}=t+a$ from \eqr{7.1}.
Sections 7.3 and 7.4 then computed what every kinematic and mechanical
variable does under such a change, and those computations are the
dictionary this section uses. The entries needed here are

\[
\begin{aligned}
   \theta^{+}&=\theta,\qquad \eps^{+}=\eps,\qquad \eta^{+}=\eta,
     &&\text{\eqr{7.37}: absolute scalars,}\\
   \bF^{+}_{\kappa^{+}}&=\bQ\bF_{\kappa}\bK^{t},
     &&\text{\eqr{7.17},}\\
   \bL^{+}&=\bQ\bL\bQ^{t}+\bOm,\qquad \bOm=\dot\bQ\bQ^{t},
     &&\text{\eqr{7.33},}\\
   \bT^{+}&=\bQ\bT\bQ^{t},
     &&\text{\eqr{7.40},}\\
   \bq^{+}&=\bQ\bq,
     &&\text{\eqr{7.44},}
\end{aligned}
\]
where $\bK$ is the fixed orthogonal two-point tensor of \eqr{7.14}
relating the two observers' reference placements, $\bX^{+}=\bK\bX+
\bc^{+}$. One entry of the list \eqr{8.7} was not computed in Chapter 7
and takes one line here. Temperature is an absolute scalar, so the
increment $\dl\theta$ across a given pair of neighbouring places is the
same for both observers, while $\dl\bx^{+}=\bQ\,\dl\bx$ by \eqr{7.3};
therefore
\[
\begin{aligned}
   \ip{\grad^{+}\theta^{+}}{\bQ\,\dl\bx}
   &=\dl\theta^{+}=\dl\theta
   \\
   &=\ip{\grad\theta}{\dl\bx}
   \quad\Longrightarrow\quad
   \ip{\big(\bQ^{t}\grad^{+}\theta^{+}-\grad\theta\big)}{\dl\bx}=0
\end{aligned}
\]
for every $\dl\bx$, whence $\bQ^{t}\grad^{+}\theta^{+}=\grad\theta$ and,
multiplying by $\bQ$ and using $\bQ\bQ^{t}=\I$,
\[
   \grad^{+}\theta^{+}=\bQ\grad\theta .
\]
The move is \eqr{8.5}--\eqr{8.6} once more, with $\bQ$ in place of
$\bF$; it is the third time in this chapter that the arbitrariness of a
line element has converted a statement about increments into a statement
about gradients.

\subsection*{What both observers must be given}

Material frame indifference is the requirement that all observers agree
on the nature of the material at hand, and hence on the class of
processes to which the material responds. With reference to \eqr{8.4}
and \eqr{8.7}, then, the response functions in the frame of $O^{+}$
should deliver the values of the elements of the set
\begin{equation}
   C^{+}=\{\bT^{+},\bq^{+},\psi^{+},\eta^{+}\}
   \tag{8.27}\label{eq:8.27}
\end{equation}
in terms of the elements of the list
\begin{equation}
   L^{+}=\{\bx^{+},\bv^{+},\theta^{+},\bF^{+}_{\kappa^{+}},\bL^{+},
   \grad^{+}\theta^{+}\},
   \tag{8.28}\label{eq:8.28}
\end{equation}
all evaluated at the material point $p\in B$ and the time $t^{+}$. So if
$O$ adopts the constitutive function
$F_{\kappa}(\bx,\bv,\theta,\bF_{\kappa},\bL,\grad\theta;\bX)$ to
determine the values of $F(p,t)\in C$, then $O^{+}$ should adopt a
function $F^{+}_{\kappa^{+}}$ of the starred arguments to evaluate its
counterpart $F^{+}(p,t^{+})\in C^{+}$.

Note what \eqr{8.27} and \eqr{8.28} are asserting, since at first sight
they assert nothing. They say that the \emph{form} of the constitutive
hypothesis is the same for both observers: $O^{+}$ has a set of four
responses determined by a list of six process variables, in one-to-one
correspondence with $O$'s. If this failed, the two observers would not be
describing the same material at all, and there would be nothing left to
compare. It is the counterpart, for a change of observer, of the
statement made after \eqr{8.12} for a change of reference.

\begin{remark}[the star on $\kappa$]\label{rem:ch8starkappa}
The reference configuration carries a star too, and it is worth saying
why this changes little. A reference configuration is a placement of the
body chosen for bookkeeping, not a thing either observer watches, and by
\eqr{7.14} the two placements differ by $\bX^{+}=\bK\bX+\bc^{+}$ with
$\bK$ a \emph{fixed orthogonal} two-point tensor and $\bc^{+}$ a fixed
vector. Fixed, because it is built in \eqr{7.15} from $\bQ(t_{0})$ and
the shifters at the initial instant, and none of those depends on $t$.

Two things follow. First, $\bK$ is a change of reference configuration
of exactly the kind treated in Section~\ref{sec:ch8noll}, so its effect
on the symmetry group is Noll's rule, and since $\bK$ is orthogonal the
conjugation is $\bK\mathcal{G}\bK^{t}$, which by the calculation after
\eqr{8.26} carries $\Orth^{+}$ into itself: a change of observer never
distorts an undistorted configuration. Second, if the two observers
agree to use the same reference placement, as one always may by taking
$\bQ(t_{0})=\I$ and equal shifters, then $\bK=\hI$, $\bc^{+}=\bze$,
$\bX^{+}=\bX$, $\kappa^{+}=\kappa$ and \eqr{7.17} reduces to
$\bF^{+}=\bQ\bF$. That is the convention used in every calculation
below. It costs nothing, and it removes a star that would otherwise
clutter every formula without changing any conclusion.
\end{remark}

\subsection*{The four transformation statements}

Take the four members of $C$ in turn. If the free energy and entropy are
determined by $O$ to be
\begin{equation}
\begin{aligned}
   \psi&=\hat\psi_{\kappa}(\bx,\bv,\theta,\bF_{\kappa},\bL,
     \grad\theta;\bX)\\
   \text{and}\quad
   \eta&=\hat\eta_{\kappa}(\bx,\bv,\theta,\bF_{\kappa},\bL,
     \grad\theta;\bX),
\end{aligned}
   \tag{8.29}\label{eq:8.29}
\end{equation}
in terms of constitutive functions $\hat\psi_{\kappa}$ and
$\hat\eta_{\kappa}$, then the corresponding constitutive equations in
the frame of $O^{+}$ are
\begin{equation}
\begin{aligned}
   \hat\psi^{+}_{\kappa^{+}}\big(\bx^{+},\bv^{+},\theta^{+},
     \bF^{+}_{\kappa^{+}},\bL^{+},\grad^{+}\theta^{+};\bX^{+}\big)
   &=\hat\psi_{\kappa}(\bx,\bv,\theta,\bF_{\kappa},\bL,
     \grad\theta;\bX)\\
   \text{and}\quad
   \hat\eta^{+}_{\kappa^{+}}\big(\bx^{+},\bv^{+},\theta^{+},
     \bF^{+}_{\kappa^{+}},\bL^{+},\grad^{+}\theta^{+};\bX^{+}\big)
   &=\hat\eta_{\kappa}(\bx,\bv,\theta,\bF_{\kappa},\bL,
     \grad\theta;\bX),
\end{aligned}
   \tag{8.30}\label{eq:8.30}
\end{equation}
respectively. There is no $\bQ$ anywhere in \eqr{8.30}, and the reason
is \eqr{6.186} together with \eqr{7.37}: the energy $\eps$, the
temperature $\theta$ and the entropy $\eta$ are absolute scalars, hence
$\psi^{+}=\eps^{+}-\theta^{+}\eta^{+}=\eps-\theta\eta=\psi$, and a
scalar has no components to rotate.

Similarly, if the heat flux is determined by $O$ to be
\begin{equation}
   \bq=\hat\bq_{\kappa}(\bx,\bv,\theta,\bF_{\kappa},\bL,\grad\theta;\bX)
   \tag{8.31}\label{eq:8.31}
\end{equation}
in terms of a constitutive function $\hat\bq_{\kappa}$, then, with
reference to \eqr{7.44}, the corresponding constitutive equation in the
frame of $O^{+}$ is
\begin{equation}
   \hat\bq^{+}_{\kappa^{+}}\big(\bx^{+},\bv^{+},\theta^{+},
     \bF^{+}_{\kappa^{+}},\bL^{+},\grad^{+}\theta^{+};\bX^{+}\big)
   =\bQ\,\hat\bq_{\kappa}(\bx,\bv,\theta,\bF_{\kappa},\bL,
     \grad\theta;\bX),
   \tag{8.32}\label{eq:8.32}
\end{equation}
where $\bQ$ is the orthogonal transformation of \eqr{7.3} that
characterizes the change of frame. The single factor $\bQ$ is the mark
of a spatial vector: one index, one rotation.

\begin{remark}[$\hat\bq_{\kappa}$ is not $\bq_{\kappa}$]
\label{rem:ch8qwarn}
The subscript $\kappa$ on $\hat\bq_{\kappa}$ records that the function
takes $\bF_{\kappa}$ among its arguments and therefore depends on the
reference configuration, exactly as in \eqr{8.10}. It does \emph{not}
mean that the output is referential. The output is the spatial heat flux
$\bq$, a vector in $\Vt$, and it transforms with one factor of $\bQ$.

The referential heat flux of \eqr{6.171}, $\bq_{\kappa}=(\bF^{*})^{t}
\bq=J\bF^{-1}\bq$, is a different object living in $\hVt$, and by
\eqr{7.45} it transforms with $\bK$ rather than $\bQ$. The two are
distinguished here by the caret, which marks a constitutive
\emph{function} rather than a value, and by nothing else, so the caret
has to be read. Getting them confused would put a $\bQ$ where a $\bK$
belongs and would make a frame-indifferent theory look frame dependent.
\end{remark}

Lastly, if $O$ finds the stress to be determined constitutively by
\begin{equation}
   \bT=\hat\bT_{\kappa}(\bx,\bv,\theta,\bF_{\kappa},\bL;\bX),
   \tag{8.33}\label{eq:8.33}
\end{equation}
then the stress is determined by $O^{+}$ to be
\begin{equation}
   \bT^{+}=\hat\bT^{+}_{\kappa^{+}}\big(\bx^{+},\bv^{+},\theta^{+},
     \bF^{+}_{\kappa^{+}},\bL^{+};\bX^{+}\big),
   \tag{8.34}\label{eq:8.34}
\end{equation}
where, by \eqr{7.40},
\begin{equation}
   \hat\bT^{+}_{\kappa^{+}}\big(\bx^{+},\bv^{+},\theta^{+},
     \bF^{+}_{\kappa^{+}},\bL^{+};\bX^{+}\big)
   =\bQ\,\hat\bT_{\kappa}(\bx,\bv,\theta,\bF_{\kappa},\bL;\bX)\,\bQ^{t}.
   \tag{8.35}\label{eq:8.35}
\end{equation}
Two factors of $\bQ$, one for each index, which is what
Theorem~\ref{thm:tenstrans} says a second-order spatial tensor must do.

The dependence of the stress on the temperature gradient has been
suppressed in \eqr{8.33} to \eqr{8.35}. This is in accord with classical
treatments of the present class of constitutive equations, and it is a
restriction on the model rather than a consequence of anything proved so
far; the entropy inequality is what eventually justifies it, and the
matter is taken up again in the chapters that follow.

\subsection*{Where the restriction actually comes from}

Equations \eqr{8.30}, \eqr{8.32} and \eqr{8.35} look like restrictions
and, taken by themselves, are not. Read literally, each of them
\emph{defines} the starred function on the left in terms of the unstarred
one on the right: given $\hat\bT_{\kappa}$, and given the dictionary of
Chapter 7 that expresses the starred arguments through the unstarred, the
right-hand side of \eqr{8.35} is a computable expression in the starred
arguments, and \eqr{8.35} names it $\hat\bT^{+}_{\kappa^{+}}$. Nothing
has been forbidden.

The physical principle is the further demand that the two observers use
\emph{the same function}:
\[
   \hat\psi^{+}_{\kappa^{+}}=\hat\psi_{\kappa},\qquad
   \hat\eta^{+}_{\kappa^{+}}=\hat\eta_{\kappa},\qquad
   \hat\bq^{+}_{\kappa^{+}}=\hat\bq_{\kappa},\qquad
   \hat\bT^{+}_{\kappa^{+}}=\hat\bT_{\kappa},
\]
under the convention $\kappa^{+}=\kappa$ of
Remark~\ref{rem:ch8starkappa}. This is the content of material frame
indifference: the elasticity of copper is a property of copper, so the
function that gives the stress of copper cannot depend on which
laboratory bench it was measured from. Impose it on \eqr{8.35} and the
definition turns into a functional equation,
\[
   \hat\bT_{\kappa}\big(\bx^{+},\bv^{+},\theta,\bQ\bF,
     \bQ\bL\bQ^{t}+\bOm;\bX\big)
   =\bQ\,\hat\bT_{\kappa}(\bx,\bv,\theta,\bF,\bL;\bX)\,\bQ^{t},
\]
required to hold for every change of frame, that is, for every choice of
the functions $\bQ(t)$ and $\bc(t)$. Now something is forbidden, and
what is forbidden is found by choosing $\bQ$ and $\bc$ cleverly.
The chapters that follow carry this out in full for the two sub-models,
viscous fluids and elastic solids. Two
consequences are worth extracting here, the first because
Section~\ref{sec:ch8elastic} promised it.

\begin{remark}[the constitutive functions cannot depend on $\bx$ or
$\bv$]\label{rem:ch8noxv}
Fix an instant $t$ later than the instant $t_{0}$ at which the reference
placements were fixed, and consider the change of frame with
\[
   \bQ(\tau)=\I\ \text{for all }\tau,
   \qquad
   \bc(t_{0})=\bze,\qquad
   \bc(t)=\ba,\qquad
   \dot\bc(t)=\bb ,
\]
where $\ba$ and $\bb$ are arbitrary vectors. Such a $\bc$ exists: three
conditions at two distinct instants are met by, for example, a cubic
polynomial in $\tau$, and no relation between $\ba$ and $\bb$ is
implied because $t\neq t_{0}$.

Work out what this frame change does to each entry of the list, at the
instant $t$.
\begin{itemize}\itemsep1pt
\item $\bQ(t_{0})=\I$ and equal shifters give $\bK=\hI$ and
      $\bc^{+}=\bze$ in \eqr{7.14}, so $\bX^{+}=\bX$ and
      $\kappa^{+}=\kappa$.
\item $\bx^{+}=\bQ\bx+\bc=\bx+\ba$ by \eqr{7.3}.
\item $\bOm=\dot\bQ\bQ^{t}=\Ob$ since $\bQ$ is constant, so \eqr{7.8}
      gives $\bv^{+}=\bv+\dot\bc=\bv+\bb$.
\item $\bF^{+}=\bQ\bF\bK^{t}=\bF$ and, by \eqr{7.33},
      $\bL^{+}=\bQ\bL\bQ^{t}+\bOm=\bL$.
\item $\theta^{+}=\theta$ and $\grad^{+}\theta^{+}=\bQ\grad\theta
      =\grad\theta$.
\end{itemize}
So this change of frame moves $\bx$ and $\bv$ and moves nothing else at
all. Since $\bQ=\I$, \eqr{8.35} reads $\hat\bT^{+}_{\kappa}=
\hat\bT_{\kappa}$, and frame indifference then gives
\[
   \hat\bT_{\kappa}(\bx+\ba,\bv+\bb,\theta,\bF,\bL;\bX)
   =\hat\bT_{\kappa}(\bx,\bv,\theta,\bF,\bL;\bX)
\]
for every $\ba$ and every $\bb$. A function whose value is unchanged
when its first two arguments are shifted by arbitrary amounts does not
depend on those arguments. The same argument applied to \eqr{8.30} and
\eqr{8.32}, in which $\bQ=\I$ likewise removes every factor of $\bQ$,
disposes of $\bx$ and $\bv$ in $\hat\psi_{\kappa}$,
$\hat\eta_{\kappa}$ and $\hat\bq_{\kappa}$ as well. This is the claim
announced without proof in Section~\ref{sec:ch8elastic}, and the reason
it holds is worth stating in words: $\bx$ is where the particle
happens to be and $\bv$ is how fast it happens to be going, and both are
artefacts of the observer's choice of origin and of state of motion. A
material cannot respond to either, because neither is a property of the
material.

The reduced lists are therefore
\[
   C=\{\bT,\bq,\psi,\eta\}
   \quad\text{determined by}\quad
   \{\theta,\bF_{\kappa},\bL,\grad\theta\},
\]
and it is these four arguments, not six, that the following chapters
work with.
\end{remark}

\begin{remark}[a preview: why viscous stress is written in $\bD$]
\label{rem:ch8Donly}
One more choice of $\bQ$ is instructive, and it explains a feature of
every viscous constitutive equation the reader has ever seen.

Choose $\bQ(\tau)$ with $\bQ(t_{0})=\I$ and $\bQ(t)=\I$ but
$\dot\bQ(t)=\bOm$ an arbitrary skew tensor; a rotation that starts at
the identity, wanders, and returns to the identity at time $t$ with
non-zero angular velocity will do. At the instant $t$ we then have
$\bQ=\I$, so
\[
   \bF^{+}=\bF,\qquad
   \theta^{+}=\theta,\qquad
   \grad^{+}\theta^{+}=\grad\theta,\qquad
   \bT^{+}=\bT,
\]
while \eqr{7.33} gives $\bL^{+}=\bL+\bOm$. Frame indifference applied to
\eqr{8.35}, with $\bx$ and $\bv$ already gone by
Remark~\ref{rem:ch8noxv}, therefore requires
\[
   \hat\bT_{\kappa}(\theta,\bF,\bL+\bOm;\bX)
   =\hat\bT_{\kappa}(\theta,\bF,\bL;\bX)
   \qquad\text{for every skew }\bOm .
\]
Choose $\bOm=-\bW$, the spin tensor of \eqr{4.7}, which is admissible
because $\bW$ is skew. Then $\bL+\bOm=\bD+\bW-\bW=\bD$ and
\[
   \hat\bT_{\kappa}(\theta,\bF,\bL;\bX)
   =\hat\bT_{\kappa}(\theta,\bF,\bD;\bX):
\]
the stress can depend on the velocity gradient only through its
symmetric part. The same holds for $\hat\psi_{\kappa}$,
$\hat\eta_{\kappa}$ and $\hat\bq_{\kappa}$.

This is the constitutive counterpart of Remark~\ref{rem:Wtrap}. The spin
$\bW$ measures rotation relative to the observer's frame, and an
observer who spins along with the material sees a different $\bW$;
\eqr{7.34} says so explicitly, $\bW^{+}=\bQ\bW\bQ^{t}+\bOm$, with the
inhomogeneous term $\bOm$ that $\bD^{+}=\bQ\bD\bQ^{t}$ does not have.
A material cannot respond to a quantity that an observer can change by
turning around, so $\bW$ must drop out and $\bD$ must remain. Every
constitutive law for a viscous fluid is written in $\bD$ for this
reason and for no other.
\end{remark}

\subsection*{What frame invariance leaves undone}

Frame invariance alone does not suffice to restrict the class of
admissible constitutive functions. Nothing said so far prevents a
material from generating energy out of nothing, or from conducting heat
from cold to hot; those are thermodynamic prohibitions and the
requirement that two observers agree cannot see them. We must also
ensure that the constitutive functions are compatible with the balances
of mass, energy and momentum and with the Clausius--Duhem inequality of
Section 6.7.

Of those four, three are quickly disposed of. Conservation of mass
imposes nothing on the list, since it merely determines $\rho$ from
$\rho_{\kappa}$ and $\det\bF$. The balance of energy imposes nothing
either, since the heat supply $r$ is at our disposal and may be chosen
to satisfy it whatever the other fields do; this is exactly the freedom
that the Coleman--Noll argument below exploits. The balance of linear
momentum imposes nothing beyond what \eqr{8.35} already contains, but
the balance of moment of momentum does: it requires $\bT=\bT^{t}$, hence
$\hat\bT_{\kappa}=\hat\bT_{\kappa}^{t}$ identically in its arguments,
which is six scalar conditions rather than nine.

The Clausius--Duhem inequality is the substantial one, and the following
remark sketches how it is used, both because it is the organising
principle of every constitutive chapter that follows and because it
shows why Section~\ref{sec:ch8elastic} took the trouble to trade $\eps$
for $\psi$.

\begin{remark}[how the entropy inequality bites]\label{rem:ch8cd}
The local inequality \eqr{6.187} of Chapter 6 is
\[
   -\rho\big(\dot\psi+\eta\dot\theta\big)+\ip\bT\bL
   -\theta^{-1}\ip\bq{\grad\theta}\ \geq\ 0 ,
\]
and it is required to hold in \emph{every} admissible process, not in
some. Substitute the constitutive equations in their reduced form,
\[
   \psi=\hat\psi_{\kappa}(\theta,\bF,\bL,\grad\theta;\bX)
\]
and so on, and expand $\dot\psi$ by the chain rule:
\[
   \dot\psi=\frac{\partial\hat\psi_{\kappa}}{\partial\theta}\dot\theta
   +\frac{\partial\hat\psi_{\kappa}}{\partial\bF}\cdot\dot\bF
   +\frac{\partial\hat\psi_{\kappa}}{\partial\bL}\cdot\dot\bL
   +\frac{\partial\hat\psi_{\kappa}}{\partial(\grad\theta)}
     \cdot(\grad\theta)^{\textstyle\cdot} .
\]
Collect terms:
\[
\begin{aligned}
   &-\rho\Big(\frac{\partial\hat\psi_{\kappa}}{\partial\theta}
   +\eta\Big)\dot\theta
   -\rho\frac{\partial\hat\psi_{\kappa}}{\partial\bL}\cdot\dot\bL
   -\rho\frac{\partial\hat\psi_{\kappa}}{\partial(\grad\theta)}
     \cdot(\grad\theta)^{\textstyle\cdot}\\
   &\qquad
   -\rho\frac{\partial\hat\psi_{\kappa}}{\partial\bF}\cdot\dot\bF
   +\ip\bT\bL-\theta^{-1}\ip\bq{\grad\theta}\ \geq\ 0 .
\end{aligned}
\]
Now the decisive observation. The three quantities $\dot\theta$,
$\dot\bL$ and $(\grad\theta)^{\textstyle\cdot}$ appear \emph{linearly},
and their coefficients do not contain them: the coefficients are
functions of $(\theta,\bF,\bL,\grad\theta)$ alone. But at a given
material point and instant one may prescribe a process realising any
values of $\dot\theta$, $\dot\bL$ and $(\grad\theta)^{\textstyle\cdot}$
one pleases, with $(\theta,\bF,\bL,\grad\theta)$ held at whatever values
are under consideration. A quantity of the form $\alpha z+\beta$, with
$\alpha$ and $\beta$ independent of $z$ and $z$ free to take any value
including large negative ones, is bounded below only if $\alpha=0$.
Hence each of the three coefficients vanishes:
\[
   \eta=-\frac{\partial\hat\psi_{\kappa}}{\partial\theta},
   \qquad
   \frac{\partial\hat\psi_{\kappa}}{\partial\bL}=\Ob,
   \qquad
   \frac{\partial\hat\psi_{\kappa}}{\partial(\grad\theta)}=\bze .
\]
Three results at once, and none of them was assumed. The free energy
depends on $\theta$ and $\bF$ only, so it is purely elastic; the
entropy is not an independent constitutive function at all but the
temperature derivative of the free energy; and the viscosity, which
lives in the dependence on $\bL$, has been banished from the energy and
confined to the stress.

What remains, after $\dot\bF=\bL\bF$ from \eqr{4.3} is used to write
$(\partial\hat\psi_{\kappa}/\partial\bF)\cdot\dot\bF=
\big[(\partial\hat\psi_{\kappa}/\partial\bF)\bF^{t}\big]\cdot\bL$, is
\[
   \Big(\bT-\rho\frac{\partial\hat\psi_{\kappa}}{\partial\bF}
   \bF^{t}\Big)\cdot\bL-\theta^{-1}\ip\bq{\grad\theta}\ \geq\ 0 ,
\]
the \emph{residual dissipation inequality}. It does not determine
anything, because $\bL$ and $\grad\theta$ are among the constitutive
arguments and so cannot be varied freely against $\bT$ and $\bq$. What
it does is split the stress into an equilibrium part
$\rho(\partial\hat\psi_{\kappa}/\partial\bF)\bF^{t}$, delivered by the
free energy, and a dissipative remainder that must do non-negative work,
and it forces heat to flow down the temperature gradient. Those two
statements are the whole thermodynamic content of the theory, and the
following chapters develop them in the contexts of distinct sub-models
for viscous fluids and elastic solids.

The sketch also settles a question left open at \eqr{8.33}: why the
temperature gradient could be dropped from the stress. It is not dropped
by fiat. Once $\hat\psi_{\kappa}$ is known to be independent of
$\grad\theta$, the residual inequality above must hold for
$\grad\theta=\bze$ as well, and the argument that then bounds the
$\grad\theta$ dependence of $\bT$ is the one carried out in the chapters
that follow.
\end{remark}

With that the general theory stops. Nothing more can be said about
constitutive equations without committing to a particular kind of
material, since everything assembled here holds equally of steel and of
honey: the locality assumption of Section~\ref{sec:ch8elastic}, the
functional equation \eqr{8.16} and the group it defines, Noll's rule
\eqr{8.26}, and the frame-invariance requirements \eqr{8.30},
\eqr{8.32} and \eqr{8.35}. The commitment is made along the fault line
that Section~\ref{sec:ch8symm} opened. Take $\mathcal{G}=U$ and the
material is a fluid, which is the subject of the next chapter; take
$\mathcal{G}\subseteq\Orth^{+}$ in some undistorted configuration and it
is a solid. The constraints of this chapter are what both of them are
built inside.
